\documentclass[11pt]{article}

\usepackage[margin=1in]{geometry}
\usepackage[T1]{fontenc}
\usepackage{lmodern}
\usepackage{amsmath,amssymb,amsthm}
\usepackage{mathtools}
\usepackage{enumitem}
\usepackage{microtype}
\usepackage[hidelinks]{hyperref}

\newcommand{\R}{\mathbb{R}}
\newcommand{\one}{\mathbf{1}}
\newcommand{\M}{\mathcal{M}}
\newcommand{\CC}{\mathcal{C}}
\newcommand{\UU}{\mathcal{U}}
\newcommand{\SR}{\mathcal{S}}
\newcommand{\Acc}{\mathrm{Acc}}
\newcommand{\hull}{\operatorname{hull}}
\newcommand{\Rp}{\R^n_{+}}
\newcommand{\Rpp}{\inter\R^n_{+}}
\DeclareMathOperator{\conv}{conv}
\DeclareMathOperator{\cone}{cone}
\DeclareMathOperator{\relint}{relint}
\DeclareMathOperator{\rbd}{rbd}
\DeclareMathOperator{\aff}{aff}
\DeclareMathOperator{\cl}{cl}
\DeclareMathOperator{\inter}{int}

\theoremstyle{plain}
\newtheorem{theorem}{Theorem}[section]
\newtheorem{proposition}[theorem]{Proposition}
\newtheorem{lemma}[theorem]{Lemma}
\newtheorem{corollary}[theorem]{Corollary}
\theoremstyle{definition}
\newtheorem{definition}[theorem]{Definition}
\newtheorem{example}[theorem]{Example}
\theoremstyle{remark}
\newtheorem{remark}[theorem]{Remark}

\setlist{itemsep=2pt,topsep=4pt}

\title{Universal coverage of coherent sets of desirable gambles}
\author{Jason Konek}
\date{Draft, 26 October 2026}

\begin{document}
\maketitle

\begin{abstract}
Fix a one-sided domain $U=D_{K'}$ of gambles. A model is a set of gambles in $U$, and a
ranking assigns to each world a reflexive relation on models. We define a ranking by a
single rule. At world $\omega$, a model $X$ does at least as well as a model $Y$ if the
change from $Y$ to $X$ is guided by the truth at $\omega$, or if the change strictly
improves the Brier score at $\omega$ and cannot be applied to any coherent model. The
second condition has an explicit form: the gambles that the change adds meet, in positive
measure, the smallest coherent model containing the gambles that it removes. The ranking is
strictly guided by the truth, compares models only where they differ, satisfies
finite-chain consistency and disjoint additivity, and has no nontrivial ties. No coherent
model has a weak dominator other than itself, and a model dominates another exactly when
the change between them is safe and improves the Brier score at every world. The main
geometric result is that every incoherent model $X$ is improved in Brier score, by a
positive margin at every world, by a coherent model $N$ whose additions are implied by its
removals: the smallest coherent model containing $X\setminus N$ is the smallest coherent
model containing $X$, and it contains $N$. The proof optimises within this coherent hull by
a minimax argument and then shrinks the optimiser; a contraction argument shows that the
shrinking preserves the hull. Consequently the undominated models are exactly the coherent
models on $U$, and every incoherent model is strictly dominated by a coherent one. With
finitely many worlds, the ranking covers exactly the full-dimensional credal sets $K$ with
$K'\subseteq\relint K$, which are all the credal sets that any ranking covers on $U$;
every coherent model on $U$ is undominated, including $D_K$ for lower-dimensional $K$, and
what such $K$ lack on a one-sided domain is local freedom. The results extend to any set
of worlds and any separable, norm-closed space of bounded gambles containing the
constants, given a measure on gambles with a finite first moment, null indifference
boundaries and full support; Lyapunov's theorem is replaced by a relaxation used only
inside the proof. There the ranking is again a universal cover on $U$, the covered credal
sets are characterised by a positivity condition, contamination models $(1-\varepsilon)
\{p\}+\varepsilon S$ with $\varepsilon>0$ on $C(\Omega)$ are covered, and, when the space
of gambles is infinite-dimensional, precise and finitely generated targets are not covered
on any one-sided domain. We isolate the argument as an abstract construction lemma. The
ranking refers to coherence in its definition: it represents the coherent models as the
undominated ones, and does not derive coherence from accuracy considerations specified
independently of coherence.
\end{abstract}

\section{Introduction}\label{sec:intro}

A model is a set of gambles, the gambles it accepts. A ranking compares models at each
world. A ranking \emph{covers} a credal set $K$ if it is guided by the truth, compares
models only where they differ, makes every undominated model coherent, is closed under
ties with undominated coherent models, makes every incoherent model dominated by a coherent
undominated one, accepts $D_K$ as a model and leaves it undominated, and lives on a domain
that leaves $D_K$ free to vary near its boundary (Definition~\ref{def:cover} lists the
exact requirements). The relations are required to be reflexive and to satisfy
the acyclicity condition (TC) below; transitivity and completeness are not required.

This note shows that, on a fixed one-sided domain $U=D_{K'}$, a single ranking makes the
undominated models exactly the coherent models on $U$ (Theorem~\ref{thm:main}). The
ranking has one rule (Definition~\ref{def:SI}). At world $i$ it accepts a comparison of
$X$ over $Y$ if the change from $Y$ to $X$ is guided by the truth at $i$, or if the change
strictly improves the Brier score at $i$ and is \emph{safe}: it cannot be applied to any
coherent model. Every accepted nontrivial comparison strictly improves one real score at
its world, which gives acyclicity at once (Proposition~\ref{prop:structure}). A change out
of a coherent model can always be applied to that model, so it is accepted only as
guidance, and guidance alone cannot weakly dominate a coherent model; this protects every
coherent model (Corollary~\ref{cor:protect}). What remains is to show that every
incoherent model is improved, at every world, by a safe change to a coherent model.

This is Theorem~\ref{thm:H}, and its proof has two steps. Let $X$ be incoherent, and let
$\hull(X)$ be the smallest coherent model containing $X$. First, a minimax argument over
the Brier moments of the subsets of $\hull(X)$ yields a coherent model $N_0\subseteq
\hull(X)$ that improves on $X$ by a positive margin at every world. Second, $N_0$ is
shrunk slightly, to a coherent model $N$, so that the gambles $X\setminus N$ that the
change from $X$ to $N$ removes still generate all of $\hull(X)$:
\[
  \hull(X\setminus N)=\hull(X)\supseteq N .
\]
This identity is proved by a contraction argument (Lemma~\ref{lem:contraction} and
Proposition~\ref{prop:shrink}). It says that every gamble the change adds is already
implied, by coherence, by the gambles the change removes: a coherent model that accepts
the removed gambles accepts the added ones too. So the change cannot be applied to any
coherent model, and the ranking accepts it, strictly, at every world.
Section~\ref{sec:abstract} separates this argument into an abstract construction lemma and
the geometric repair theorem.

Section~\ref{sec:infinite} extends the results to an arbitrary set of worlds $\Omega$ and a
separable, norm-closed space $E$ of bounded gambles containing the constants. The measure
now lives on $E$; it needs a finite first moment, null indifference boundaries and, for
the coverage results, full support, and Lemma~\ref{lem:measure} constructs one. Lyapunov's theorem is replaced by a relaxation:
acceptance sets are temporarily allowed to accept gambles with a weight between $0$ and
$1$, the minimax theorem is applied to the relaxed choices, and uniqueness of the optimal
choice returns an ordinary model (Theorem~\ref{thm:Hinf}). The contraction argument is
unchanged. The safe-improvement ranking again makes the undominated models exactly the
coherent ones (Theorem~\ref{thm:inf-main}) and is a universal cover on $U$; the covered
credal sets are characterised by a positivity condition that replaces ``full-dimensional
with $K'$ in the relative interior'' (Theorem~\ref{thm:inf-cover}). Contamination models
with $\varepsilon>0$ on $C(\Omega)$ are covered (Proposition~\ref{prop:contamination}),
while, when $E$ is infinite-dimensional, precise and finitely generated targets are not
covered on any one-sided domain. Section~\ref{sec:scope} records the scope
of the results.

\section{Setting and framework}\label{sec:framework}

\subsection{Gambles, credal sets and the measure}

Let $\Omega=\{1,\dots,n\}$ with $n\ge2$. A \emph{gamble} is a vector $h\in\R^n$; $h_i$ is
its payoff at world $i$. Write $e_1,\dots,e_n$ for the standard basis, $\one\coloneqq
(1,\dots,1)$, $\Delta\coloneqq\{m\in\Rp:\one\cdot m=1\}$ for the simplex and
$\aff\Delta\coloneqq\{m:\one\cdot m=1\}$. For $C\subseteq\aff\Delta$, $\relint C$ is the interior
of $C$ relative to $\aff\Delta$ and $\rbd C\coloneqq C\setminus\relint C$. A
\emph{credal set} is a nonempty closed convex subset of $\Delta$; it is
\emph{full-dimensional} if $\relint K\neq\emptyset$. For $S\subseteq\R^n$, $\cone S$ is
the set of finite non-negative combinations of points of $S$, with $\cone\emptyset=\{0\}$.
We write $x\gg0$ if $x\in\Rpp$, that is, if every coordinate of $x$ is positive.

For $S\subseteq\Delta$ define
\[
  D_S\coloneqq\{h\in\R^n: q\cdot h>0\text{ for all }q\in S\},\qquad
  C_S\coloneqq\{h\in\R^n: q\cdot h\ge0\text{ for all }q\in S\},
\]
with $D_\emptyset=C_\emptyset\coloneqq\R^n$, and, for nonempty compact $S$,
$P_S(h)\coloneqq\min_{q\in S}q\cdot h$. Note that $D_\Delta=\Rpp$.

\begin{itemize}
  \item[(M)] $\mu$ is a finite Borel measure on $\R^n$ with the same null sets as Lebesgue
  measure, and $\int|h|\,d\mu(h)<\infty$.
\end{itemize}
Throughout, $\mu$ satisfies (M). Equalities and inclusions between sets of gambles hold
\emph{almost everywhere} (a.e.) if they hold up to a $\mu$-null set. Every hyperplane and
every sphere is null, a nonempty open set is never null, and $\mu$ is non-atomic. For
measurable $X$ put
\[
  m(X)\coloneqq\int_Xh\,d\mu(h)\in\R^n,\qquad m_i(X)\coloneqq\int_Xh_i\,d\mu(h).
\]
For measurable $E$, $E_d$ denotes the set of Lebesgue density points of $E$. By the
Lebesgue density theorem $E\setminus E_d$ is null, and by (M), $\mu(E\cap V)>0$ for every
open ball $V$ centred at a point of $E_d$.

\begin{lemma}\label{lem:D}
Let $S,S_1,S_2\subseteq\Delta$ be nonempty and compact.
\begin{enumerate}[label=(\alph*)]
  \item $D_S$ is an open convex cone containing $\Rpp$, and $D_S\cap(-D_S)=\emptyset$.
  Moreover $D_S=D_{\conv S}$ and $D_{S_1}\cap D_{S_2}=D_{S_1\cup S_2}$.
  \item $P_S(h+t\one)=P_S(h)+t$ for $t\in\R$. The boundary $\partial D_S$ is contained in
  $\{P_S=0\}$, and $\{P_S=0\}$ is null.
  \item If $S_1,S_2$ are convex and $D_{S_1}\subseteq D_{S_2}$, then $S_2\subseteq S_1$.
  \item If $O$ is open and $O\subseteq D_S$ a.e., then $O\subseteq D_S$.
\end{enumerate}
\end{lemma}

\begin{proof}
(a) $|P_S(h)-P_S(h')|\le\max_j|h_j-h'_j|$ because $S\subseteq\Delta$, so $D_S=\{P_S>0\}$
is open. Convexity, the cone property and $\Rpp\subseteq D_S$ are immediate, and $q\cdot
h>0$, $q\cdot(-h)>0$ cannot both hold. A linear function has the same minimum over $S$
and over $\conv S$. The last identity is the definition.
(b) The identity holds because $q\cdot\one=1$ for $q\in\Delta$. If $P_S(h)<0$ then $h$
has a neighbourhood outside $D_S$; if $P_S(h)=0$ then $h+t\one\in D_S$ for $t>0$ and
$h\notin D_S$. So $\partial D_S\subseteq\{P_S=0\}$. Write $h=y+t\one$ with $y\in
\one^\perp$, a linear change of coordinates. Then $\{P_S=0\}=\{y-P_S(y)\one:y\in
\one^\perp\}$ meets each line $y+\R\one$ in one point, so it is Lebesgue-null by Fubini's
theorem, and $\mu$-null by (M).
(c) Suppose $k\in S_2\setminus S_1$. Strict separation of $k$ from $S_1$ gives $g\in\R^n$
and $c\in\R$ with $g\cdot q>c>g\cdot k$ for all $q\in S_1$. Put $h\coloneqq g-c\one$. Then
$q\cdot h>0$ on $S_1$ and $k\cdot h<0$, so $h\in D_{S_1}\setminus D_{S_2}$.
(d) The open set $O\cap\{P_S<0\}$ is null, hence empty. If $h\in O$, then $h-\varepsilon
\one\in O$ for some $\varepsilon>0$, so $P_S(h)-\varepsilon\ge0$ by (b), and $h\in D_S$.
\end{proof}

\subsection{Models, coherence and the coherent hull}

A \emph{domain} is a measurable set $U\subseteq\R^n$. The \emph{models} on $U$ are the
measurable sets $X\subseteq U$, identified when they agree a.e.; their class is $\M_U$.
A model $X$ \emph{accepts} $h$ if $h\in X$. At world $i$, accepting $h$ is \emph{correct}
if $h_i>0$ and \emph{wrong} if $h_i<0$. A model $X$ is \emph{coherent} if $X=D_Q$ a.e.\
for some credal set $Q$. If $D_{Q_1}=D_{Q_2}$ a.e., Lemma~\ref{lem:D}(d) gives $D_{Q_1}=
D_{Q_2}$, and Lemma~\ref{lem:D}(c) gives $Q_1=Q_2$; so $Q$ is unique. We call $U$
\emph{one-sided} if $U\cap(-U)$ is null. This usage differs from that of \cite{KCM},
where ``one-sided'' means $U\cap(-U)=\emptyset$ and ``essentially one-sided'' means that
$U\cap(-U)$ contains no non-null set of gambles with a positive and a negative coordinate;
our condition lies between the two.

\smallskip\noindent\textbf{Standing assumption.} Except where a general domain is
mentioned explicitly, $K'\subseteq\Delta$ is nonempty, compact and convex, and
$U\coloneqq D_{K'}$. By Lemma~\ref{lem:D}(a), $U$ is open, $U\supseteq\Rpp$ and $U\cap
(-U)=\emptyset$; moreover $U$ contains no $h$ with $h\le0$, since $k\cdot h\le0$ for
$k\in K'$.

\smallskip
Let $\CC_U$ be the class of coherent models on $U$. If $Q$ is a credal set with
$K'\subseteq Q$, then $D_Q\subseteq D_{K'}=U$, so $D_Q\in\CC_U$. Conversely, if $D_Q
\subseteq U$ a.e., then $D_Q\subseteq D_{K'}$ by Lemma~\ref{lem:D}(d), and $K'\subseteq
Q$ by Lemma~\ref{lem:D}(c). Hence
\[
  \CC_U=\{D_Q: Q\text{ a credal set with }K'\subseteq Q\}.
\]
Every coherent model contains $\Rpp$. For measurable $B\subseteq U$ put
\[
  Q_B\coloneqq\{q\in\Delta: q\cdot h\ge0\text{ for $\mu$-a.e.\ }h\in B\},\qquad
  \hull(B)\coloneqq D_{Q_B}.
\]
If $B$ is null, then $Q_B=\Delta$ and $\hull(B)=\Rpp$.

\begin{lemma}\label{lem:hull}
Let $B\subseteq U$ be measurable.
\begin{enumerate}[label=(\alph*)]
  \item $Q_B$ is a credal set and $K'\subseteq Q_B$. Hence $\hull(B)\in\CC_U$.
  \item $B\subseteq\hull(B)$ a.e.
  \item For a credal set $Q$ with $K'\subseteq Q$, the following are equivalent:
  $B\subseteq D_Q$ a.e.; $Q\subseteq Q_B$; $\hull(B)\subseteq D_Q$. So $\hull(B)$ is the
  smallest coherent model containing $B$ a.e.
\end{enumerate}
\end{lemma}

\begin{proof}
(a) $K'\subseteq Q_B$ because $k\cdot h>0$ on $U$ for $k\in K'$. Convexity is clear. Let
$q_k\in Q_B$ converge to $q$. If $q\cdot h<0$ then $q_k\cdot h<0$ for all large $k$, so
$B\cap\{q\cdot h<0\}\subseteq\bigcup_kB\cap\{q_k\cdot h<0\}$, which is null. So $Q_B$ is
closed. Since $Q_B\supseteq K'$, $\hull(B)\in\CC_U$.
(b) $B\setminus D_{Q_B}\subseteq\{P_{Q_B}=0\}\cup\bigcup_{q\in Q_B}\{q\cdot h<0\}$. The
first set is null by Lemma~\ref{lem:D}(b). The open sets $\{q\cdot h<0\}$ have a countable
subcover $\{q_k\cdot h<0\}$ (Lindel\"of), and each $B\cap\{q_k\cdot h<0\}$ is null.
(c) If $B\subseteq D_Q$ a.e., every $q\in Q$ has $q\cdot h>0$ for a.e.\ $h\in B$, so
$Q\subseteq Q_B$. If $Q\subseteq Q_B$, then $D_{Q_B}\subseteq D_Q$ by definition. If
$\hull(B)\subseteq D_Q$, then $B\subseteq D_Q$ a.e.\ by (b).
\end{proof}

\subsection{Rankings and dominance}

A \emph{ranking} on $\M_U$ is a family $(\preceq_i)_{i\in\Omega}$ of binary relations on
$\M_U$. Lower is better: $X\preceq_iY$ says that $X$ does at least as well as $Y$ at
world $i$. We write $X\prec_iY$ if $X\preceq_iY$ and not $Y\preceq_iX$; $X$ and $Y$ are
\emph{tied} at $i$ if $X\preceq_iY$ and $Y\preceq_iX$, and \emph{incomparable} at $i$ if
neither holds. We assume throughout that each $\preceq_i$ is reflexive, that a.e.-equal
models are tied, and
\begin{itemize}
  \item[(TC)] if $A\preceq_iB$ and $B\preceq_iC$, then not $C\prec_iA$.
\end{itemize}
Transitivity and completeness are not assumed. We also consider the stronger condition
\begin{itemize}
  \item[(FCC)] \emph{Finite-chain consistency.} If $A_0\preceq_iA_1\preceq_i\dots
  \preceq_iA_k$, then not $A_k\prec_iA_0$.
\end{itemize}
(TC) is the case $k=2$ of (FCC).

\begin{definition}\label{def:dom}
Let $X,Y\in\M_U$. $Y$ \emph{weakly dominates} $X$ if $Y\preceq_iX$ for every $i$. $Y$
\emph{dominates} $X$ if $Y$ weakly dominates $X$ and $Y\prec_iX$ for some $i$. $Y$
\emph{strictly dominates} $X$ if $Y\prec_iX$ for every $i$. $X$ is \emph{undominated} if
no model dominates it. We write $\UU$ for the undominated models, $\CC$ for the coherent
models and $\CC^\ast\coloneqq\CC\cap\UU$.
\end{definition}

The \emph{Brier loss} of $X\in\M_U$ at world $i$ is the $|h_i|$-weighted measure of its
wrong acceptances and of its silences on correct gambles:
\[
  \beta_i(X)\coloneqq\int_{X\cap\{h_i<0\}}|h_i|\,d\mu
  +\int_{(U\setminus X)\cap\{h_i>0\}}|h_i|\,d\mu=c_i-m_i(X),\qquad
  c_i\coloneqq\int_{U\cap\{h_i>0\}}h_i\,d\mu .
\]
The \emph{Brier ranking} is $X\preceq^{\mathrm B}_iY$ iff $\beta_i(X)\le\beta_i(Y)$, that
is, iff $m_i(X)\ge m_i(Y)$. $Y$ \emph{uniformly dominates} $X$ in Brier loss if there is
$\varepsilon>0$ with $\beta_i(Y)+\varepsilon\le\beta_i(X)$ for every $i$, that is, if
$m(Y)-m(X)\gg0$.

\subsection{Requirements}

\begin{itemize}
  \item[(G)] \emph{Strict guidance.} If $X\setminus Y\subseteq\{h_i>0\}$ a.e.,
  $Y\setminus X\subseteq\{h_i<0\}$ a.e., and $X\triangle Y$ is not null, then
  $X\prec_iY$.
  \item[(IB)] \emph{Independence.} Let $S$ be measurable and $X,Y,X',Y'\in\M_U$ with
  $X\setminus S=Y\setminus S$, $X'\setminus S=Y'\setminus S$, $X\cap S=X'\cap S$ and
  $Y\cap S=Y'\cap S$, all a.e. Then $X\preceq_iY$ if and only if $X'\preceq_iY'$.
  \item[(GC)] \emph{Global coherence.} Every undominated model is coherent.
  \item[(C)] \emph{Direct closure under ties.} A model that is tied at every world with a
  member of $\CC^\ast$ belongs to $\CC^\ast$.
  \item[(ID)] \emph{Incoherent-model dominance.} Every incoherent model is dominated by a
  member of $\CC^\ast$.
  \item[(LF)] \emph{Local freedom at $D_K$.} There is an open set $N\supseteq\partial
  D_K\setminus\{0\}$ such that every measurable $X$ with $X\triangle D_K\subseteq N$ a.e.\
  belongs to $\M_U$, that is, satisfies $X\subseteq U$ a.e.
  \item[(W)] Every model is weakly dominated by a member of $\CC^\ast$.
\end{itemize}
We also consider
\begin{itemize}
  \item[(Add)] \emph{Disjoint additivity.} If $X\preceq_iY$ and $Y\preceq_iZ$, and
  $(X\triangle Y)\cap(Y\triangle Z)$ is null, then $X\preceq_iZ$.
\end{itemize}
(Add) is the case of transitivity in which the two steps change disjoint sets of gambles.

\begin{definition}\label{def:cover}
Let $K$ be a credal set and $U$ a domain. A ranking on $\M_U$ \emph{covers} $K$ if:
\begin{enumerate}[label=(\roman*)]
  \item it satisfies (G) and (IB);
  \item it satisfies (GC);
  \item it satisfies (C);
  \item $D_K\in\M_U$ and $D_K$ is undominated;
  \item it satisfies (ID);
  \item (LF) holds at $D_K$.
\end{enumerate}
\end{definition}

(W) and (Add) are not part of the definition; the ranking constructed below satisfies
both.

\begin{lemma}\label{lem:LF}
Let $K$ be a credal set.
\begin{enumerate}[label=(\alph*)]
  \item $D_K\in\M_U$ if and only if $K'\subseteq K$.
  \item If $K$ is full-dimensional and $K'\subseteq K$, then (LF) holds at $D_K$ if and
  only if $K'\subseteq\relint K$.
  \item Let $V$ be any one-sided domain. If $K$ is not full-dimensional, then (LF) fails
  at $D_K$ for $\M_V$.
\end{enumerate}
\end{lemma}

\begin{proof}
(a) This is the description of $\CC_U$ above.

(b) First let $h\in C_K\setminus\{0\}$ and $x\in\relint K$. If $x\cdot h=0$, then for
every $d$ with $\one\cdot d=0$ and small $t>0$ both $x\pm td$ lie in $K$, so $\pm t\,d
\cdot h\ge0$ and $d\cdot h=0$. Then $h=c\one$ for some $c$, and $c=x\cdot h=0$, a
contradiction. So $x\cdot h>0$.

Suppose $K'\subseteq\relint K$. Then $D_K\subseteq U$. Each $h\in\partial D_K\setminus
\{0\}$ lies in $C_K\setminus\{0\}$, so $k\cdot h>0$ for $k\in K'$ by the previous
paragraph, that is, $h\in U$. Take $N\coloneqq U$, which is open. If $X\triangle D_K
\subseteq U$ a.e., then $X\subseteq D_K\cup U=U$ a.e.

Conversely, suppose (LF) holds with the open set $N$, and $x\in K'\cap\rbd K$. Since
$x\notin\relint K$, which is the relative interior of $K$ in the sense of convex analysis
because $K$ is full-dimensional, there is a nontrivial supporting hyperplane to $K$
through $x$ \cite[Thm.~11.6]{Rockafellar}: some $g\in\R^n$ has $g\cdot k\ge g\cdot x$ for
all $k\in K$, with strict inequality for some $k\in K$. Put $h\coloneqq g-(g\cdot x)\one$.
Then $k\cdot h\ge0$ on $K$, $x\cdot h=0$ and $h\ne0$. So $h\in C_K\setminus D_K$, and
$h+t\one\in D_K$ for $t>0$; hence $h\in\partial D_K\setminus\{0\}\subseteq N$. The model
$D_K\cup N$ satisfies $(D_K\cup N)\triangle D_K\subseteq N$, so it belongs to $\M_U$;
thus $N\subseteq U$ a.e., and $N\subseteq U$ by Lemma~\ref{lem:D}(d). But $x\cdot h=0$
with $x\in K'$, so $h\notin U$, a contradiction.

(c) The linear span of $K$ has dimension at most $n-1$, so some $h\ne0$ has $k\cdot h=0$
for all $k\in K$. Then $\pm h\in C_K\setminus D_K$ and $\pm h+t\one\in D_K$ for $t>0$, so
$\pm h\in\partial D_K\setminus\{0\}$. If (LF) held with the open set $N$, then, as in (b),
$N\subseteq V$ a.e.; choose $\rho>0$ with $\operatorname{Ball}(\pm h,\rho)\subseteq N$,
where $\operatorname{Ball}(h,\rho)$ is the open ball of radius $\rho$ about $h$. Then
$\operatorname{Ball}(h,\rho)\subseteq V$ a.e.\ and $\operatorname{Ball}(h,\rho)=
-\operatorname{Ball}(-h,\rho)\subseteq-V$ a.e., because Lebesgue-null sets are symmetric;
so the non-null ball $\operatorname{Ball}(h,\rho)$ is contained
in $V\cap(-V)$ a.e., which contradicts one-sidedness.
\end{proof}

\section{The safe-improvement ranking}\label{sec:ranking}

\subsection{Records and safety}

A \emph{record} is an ordered pair $r=(A,B)$ of disjoint measurable subsets of $U$, with
records identified when their components agree a.e. The record of a pair of models is
\[
  r(X,Y)\coloneqq(X\setminus Y,\;Y\setminus X).
\]
It describes the change from $Y$ to $X$: the change adds $A$ and removes $B$. The record
$r$ is \emph{null} if $A$ and $B$ are null; $r(X,Y)$ is null if and only if $X=Y$ a.e.
The \emph{support} of $r$ is $A\cup B$, and its \emph{moment} is $m(r)\coloneqq
m(A)-m(B)$, so that
\[
  m_i(r(X,Y))=m_i(X)-m_i(Y)=\beta_i(Y)-\beta_i(X).
\]
For records $r_1=(A_1,B_1)$ and $r_2=(A_2,B_2)$ whose supports meet in a null set, put
$r_1\sqcup r_2\coloneqq(A_1\cup A_2,B_1\cup B_2)$; then $m(r_1\sqcup r_2)=m(r_1)+m(r_2)$.

A record $r=(A,B)$ \emph{applies to} a model $Z$ if $B\subseteq Z$ a.e.\ and $A\cap Z$ is
null: the change it describes can be carried out on $Z$. A record is \emph{safe} if it
applies to no coherent model on $U$.

\begin{lemma}\label{lem:record}
\begin{enumerate}[label=(\alph*)]
  \item A record $r$ applies to $Z\in\M_U$ if and only if $r=r(Y,Z)$ for some
  $Y\in\M_U$. In particular $r(Y,C)$ is not safe for any $Y\in\M_U$ and $C\in\CC_U$.
  \item If $X,Y,Z\in\M_U$ and $(X\triangle Y)\cap(Y\triangle Z)$ is null, then the
  supports of $r(X,Y)$ and $r(Y,Z)$ meet in a null set and $r(X,Z)=r(X,Y)\sqcup r(Y,Z)$.
  \item Let $\Acc$ be a set of records containing the null record, and define $X\preceq Y$
  iff $r(X,Y)\in\Acc$. Then $\preceq$ is reflexive, ties a.e.-equal models and satisfies
  (IB).
\end{enumerate}
\end{lemma}

\begin{proof}
(a) Let $r=(A,B)$ apply to $Z$, and put $Y\coloneqq(Z\setminus B)\cup A\in\M_U$. Then
$Y\setminus Z=A\setminus Z=A$ a.e., because $A\cap Z$ is null, and $Z\setminus Y=(Z\cap
B)\setminus A=B$ a.e., because $B\subseteq Z$ a.e.\ and $A\cap B=\emptyset$. So
$r(Y,Z)=r$. Conversely, $r(Y,Z)=(Y\setminus Z,Z\setminus Y)$ applies to $Z$ by
definition.

(b) The supports are $X\triangle Y$ and $Y\triangle Z$. Outside the null set $(X\triangle
Y)\cap(Y\triangle Z)$, every $h$ lies outside $X\triangle Y$ or outside $Y\triangle Z$.
In the first case $h\in X\iff h\in Y$, so $h\notin X\setminus Y$, and $h\in X\setminus Z$
iff $h\in Y\setminus Z$. In the second case $h\in Y\iff h\in Z$, so $h\notin Y\setminus
Z$, and $h\in X\setminus Z$ iff $h\in X\setminus Y$. In both cases $\one_{X\setminus Z}(h)
=\one_{X\setminus Y}(h)+\one_{Y\setminus Z}(h)$, and the same argument gives
$\one_{Z\setminus X}(h)=\one_{Y\setminus X}(h)+\one_{Z\setminus Y}(h)$.

(c) If $X=Y$ a.e., then $r(X,Y)$ and $r(Y,X)$ are null. For (IB), the hypotheses give
$X\triangle Y\subseteq S$ a.e., so $X\setminus Y=(X\cap S)\setminus(Y\cap S)=(X'\cap S)
\setminus(Y'\cap S)=X'\setminus Y'$ a.e., and likewise $Y\setminus X=Y'\setminus X'$
a.e. So $r(X,Y)=r(X',Y')$.
\end{proof}

\begin{lemma}\label{lem:safe}
Let $r=(A,B)$ be a record.
\begin{enumerate}[label=(\alph*)]
  \item $r$ is safe if and only if $\mu\bigl(A\cap\hull(B)\bigr)>0$.
  \item The null record is not safe. If $B$ is null, then $r$ is safe if and only if
  $\mu(A\cap\Rpp)>0$.
  \item If $r_1$ is safe and the supports of $r_1$ and $r_2$ meet in a null set, then
  $r_1\sqcup r_2$ is safe.
\end{enumerate}
\end{lemma}

\begin{proof}
(a) If $A\cap\hull(B)$ is null, then $r$ applies to $\hull(B)$, which is a coherent model
on $U$ containing $B$ a.e.\ (Lemma~\ref{lem:hull}(a),(b)); so $r$ is not safe.
Conversely, if $r$ applies to a coherent model $D_Q\in\CC_U$, then $B\subseteq D_Q$ a.e.,
so $\hull(B)\subseteq D_Q$ by Lemma~\ref{lem:hull}(c), and $A\cap\hull(B)\subseteq A\cap
D_Q$ is null.

(b) If $A$ is null, so is $A\cap\hull(B)$. If $B$ is null, then $\hull(B)=\Rpp$.

(c) Write $r_k=(A_k,B_k)$. If $r_1\sqcup r_2$ applied to a model $Z$, then $B_1\subseteq
B_1\cup B_2\subseteq Z$ a.e.\ and $A_1\cap Z\subseteq(A_1\cup A_2)\cap Z$ is null, so
$r_1$ would apply to $Z$.
\end{proof}

By (a), safety of $r(X,Y)$ says that some of the gambles the change from $Y$ to $X$ adds
are required, by coherence, of the gambles it removes. Part~(c) says that safety cannot be
lost by adjoining a further, disjoint change.

\subsection{The ranking}

For $i\in\Omega$ let
\[
  G_i\coloneqq\{r=(A,B): A\subseteq\{h_i>0\}\text{ and }B\subseteq\{h_i<0\}\text{
  a.e.}\}
\]
be the set of records that are \emph{guidance at $i$}: the change adds only gambles that
are correct at $i$ and removes only gambles that are wrong at $i$. The null record is
guidance at every world.

\begin{definition}\label{def:SI}
The \emph{safe-improvement ranking} on $\M_U$ is given by
\[
  \Acc_i\coloneqq G_i\cup\{r: r\text{ is safe and }m_i(r)>0\},\qquad
  X\preceq_iY\iff r(X,Y)\in\Acc_i .
\]
\end{definition}

In words: at world $i$, $X$ does at least as well as $Y$ if the change from $Y$ to $X$ is
guidance at $i$, or if it is safe and strictly reduces the Brier loss at $i$. No
lexicographic tie-breaking, aggregation across worlds or designated world is involved:
whether a comparison is accepted at $i$ depends only on its record, through guidance at
$i$, the Brier increment $m_i$ and the world-independent safety test.

\begin{lemma}\label{lem:increment}
\begin{enumerate}[label=(\alph*)]
  \item If $r=(A,B)\in G_i$, then $m_i(r)=\int_A|h_i|\,d\mu+\int_B|h_i|\,d\mu\ge0$, with
  equality if and only if $r$ is null.
  \item If $r\in\Acc_i$, then $m_i(r)\ge0$, with equality if and only if $r$ is null.
  \item A record lies in $G_i$ for every $i$ if and only if it is a.e.\ of the form
  $(A,\emptyset)$ with $A\subseteq\Rpp$. Every non-null such record is safe.
\end{enumerate}
\end{lemma}

\begin{proof}
(a) The formula holds because $h_i>0$ a.e.\ on $A$ and $h_i<0$ a.e.\ on $B$. The integral
vanishes iff $A\cup B\subseteq\{h_i=0\}$ a.e., and this hyperplane is null by (M).
(b) Records in $G_i$ are covered by (a); the others have $m_i(r)>0$ and are non-null.
(c) If $r=(A,B)\in G_i$ for every $i$, then $A\subseteq\bigcap_i\{h_i>0\}=\Rpp$ a.e.\
and $B\subseteq\{h:h_i<0\text{ for all }i\}$ a.e. By the standing assumption $U$
contains no $h\le0$, so $B$ is null. The converse is clear. If $A\subseteq\Rpp$ is
non-null, then $\mu(A\cap\Rpp)>0$, and $(A,\emptyset)$ is safe by
Lemma~\ref{lem:safe}(b).
\end{proof}

\begin{proposition}\label{prop:structure}
Consider the safe-improvement ranking on $\M_U$.
\begin{enumerate}[label=(\alph*)]
  \item Each $\preceq_i$ is reflexive, ties a.e.-equal models and satisfies (IB).
  \item If $X\preceq_iY$, then $m_i(X)\ge m_i(Y)$, with strict inequality unless $X=Y$
  a.e. Consequently two models are tied at $i$ only if they are a.e.\ equal, and
  $X\prec_iY$ if and only if $X\preceq_iY$ and $X\ne Y$ a.e. Each $\preceq_i$ is contained
  in the Brier ranking $\preceq^{\mathrm B}_i$, and each $\prec_i$ in its strict part.
  \item (G) holds.
  \item (FCC) holds, and hence (TC). More strongly, if $A_0\preceq_iA_1\preceq_i\dots
  \preceq_iA_k\preceq_iA_0$, then $A_0,\dots,A_k$ are a.e.\ equal.
  \item (Add) holds.
\end{enumerate}
\end{proposition}

\begin{proof}
(a) The null record lies in $G_i\subseteq\Acc_i$; apply Lemma~\ref{lem:record}(c).

(b) $m_i(X)-m_i(Y)=m_i(r(X,Y))$, and $r(X,Y)$ is null iff $X=Y$ a.e.; apply
Lemma~\ref{lem:increment}(b). If $X\preceq_iY$ and $Y\preceq_iX$ with $X\ne Y$ a.e., then
$m_i(X)>m_i(Y)>m_i(X)$, which is impossible. By (a), a.e.-equal models are tied, so
$X\prec_iY$ iff $X\preceq_iY$ and $X\ne Y$ a.e.

(c) If the hypotheses of (G) hold, then $r(X,Y)\in G_i\subseteq\Acc_i$ and $X\ne Y$
a.e.; by (b), $X\prec_iY$.

(d) By (b), a chain $A_0\preceq_i\dots\preceq_iA_k$ gives $m_i(A_0)\ge m_i(A_k)$. If
$A_k\prec_iA_0$, then $A_k\ne A_0$ a.e.\ and (b) gives $m_i(A_k)>m_i(A_0)$, a
contradiction. For a cycle, (b) gives $m_i(A_0)\ge m_i(A_1)\ge\dots\ge m_i(A_k)\ge
m_i(A_0)$, so all these inequalities are equalities, and each step relates a.e.-equal
models by (b).

(e) Let $X\preceq_iY$ and $Y\preceq_iZ$ with $(X\triangle Y)\cap(Y\triangle Z)$ null, and
put $r_1\coloneqq r(X,Y)$, $r_2\coloneqq r(Y,Z)$. By Lemma~\ref{lem:record}(b),
$r(X,Z)=r_1\sqcup r_2$. If $r_1,r_2\in G_i$, then $r_1\sqcup r_2\in G_i$, because
$A_1\cup A_2\subseteq\{h_i>0\}$ and $B_1\cup B_2\subseteq\{h_i<0\}$ a.e. Otherwise one of
them, say $r_k$, lies in $\Acc_i\setminus G_i$; then $r_k$ is safe and $m_i(r_k)>0$. By
Lemma~\ref{lem:safe}(c), $r_1\sqcup r_2$ is safe, and by Lemma~\ref{lem:increment}(b)
the other record has non-negative increment, so $m_i(r_1\sqcup r_2)=m_i(r_1)+m_i(r_2)>0$.
In both cases $r(X,Z)\in\Acc_i$, that is, $X\preceq_iZ$.
\end{proof}

\begin{proposition}\label{prop:dominance}
Consider the safe-improvement ranking on $\M_U$, and let $X,Y\in\M_U$. The following are
equivalent.
\begin{enumerate}[label=(\roman*)]
  \item $Y$ weakly dominates $X$, and $Y\ne X$ a.e.
  \item $Y$ dominates $X$.
  \item $Y$ strictly dominates $X$.
  \item $r(Y,X)$ is safe and $m(Y)-m(X)\gg0$; that is, the change from $X$ to $Y$ is safe
  and $Y$ uniformly dominates $X$ in Brier loss.
\end{enumerate}
\end{proposition}

\begin{proof}
(iii)$\Rightarrow$(ii) is immediate. (ii)$\Rightarrow$(i): if $Y\prec_jX$ for some $j$,
then $X$ and $Y$ are not tied at $j$, so $Y\ne X$ a.e.

(i)$\Rightarrow$(iv): put $r\coloneqq r(Y,X)$, which is non-null and lies in $\Acc_i$ for
every $i$. By Proposition~\ref{prop:structure}(b), $m_i(Y)>m_i(X)$ for every $i$. If $r$
were not safe, it would lie in $G_i$ for every $i$, by the definition of $\Acc_i$; but
then it would be safe by Lemma~\ref{lem:increment}(c), since it is non-null.

(iv)$\Rightarrow$(iii): $r(Y,X)$ is safe with $m_i(r(Y,X))>0$, so $Y\preceq_iX$ for every
$i$. Since $m(Y)\ne m(X)$, $Y\ne X$ a.e., and Proposition~\ref{prop:structure}(b) gives
$Y\prec_iX$ for every $i$.
\end{proof}

In particular, dominance in this ranking is never merely weak at some world: whenever a
model dominates another, it improves its Brier score at every world, by a safe change.

\begin{corollary}\label{cor:protect}
Consider the safe-improvement ranking on $\M_U$. Every coherent model on $U$ has no weak
dominator other than itself; in particular it is undominated. A model tied with a coherent
model at some world equals it a.e.
\end{corollary}

\begin{proof}
Let $C\in\CC_U$ and let $Y$ weakly dominate $C$. The record $r(Y,C)=(Y\setminus C,
C\setminus Y)$ applies to $C$, so it is not safe (Lemma~\ref{lem:record}(a)); since it is
accepted at every world, it lies in $G_i$ for every $i$. Hence $Y\setminus C\subseteq
\Rpp$ a.e., and $\Rpp\subseteq C$ a.e.\ because $C$ is coherent; so $Y\setminus C$ is
null. Also $C\setminus Y\subseteq\{h:h_i<0\text{ for all }i\}$ a.e., and no coherent model
contains such a gamble, since $q\cdot h<0$ for $q\in\Delta$; so $C\setminus Y$ is null.
Thus $Y=C$ a.e. A dominator of $C$ would be a weak dominator different from $C$. The last
claim is Proposition~\ref{prop:structure}(b).
\end{proof}

The proof uses only that every coherent model contains $\Rpp$ and that no coherent model
contains a gamble $h$ with $h_i<0$ for all $i$.

The mechanism of protection is this. A change out of a coherent model $C$ applies to
$C$, so it is never safe, and the ranking accepts it only where it is guidance. A change
that is guidance at every world only adds gambles that are positive at every world, and
these every coherent model already accepts.

\subsection{A geometric form of the safety test}

The safety test of Lemma~\ref{lem:safe}(a) has a finite geometric form.

\begin{lemma}\label{lem:hullformula}
For measurable $B\subseteq U$,
\[
  \hull(B)=\cone(B_d)+\Rpp .
\]
\end{lemma}

\begin{proof}
First, if $q\in Q_B$ and $b\in B_d$, then $q\cdot b\ge0$: otherwise an open ball $V$
around $b$ lies in $\{q\cdot h<0\}$, and $\mu(B\cap V)>0$ contradicts $q\in Q_B$.
Conversely, if $q\in\Delta$ and $q\cdot b\ge0$ for all $b\in B_d$, then $q\in Q_B$,
because a.e.\ point of $B$ lies in $B_d$.

Put $W\coloneqq\cone(B_d)+\Rpp$, a nonempty open convex set. If $a=\sum_jt_jb_j+p$
with $b_j\in B_d$, $t_j\ge0$ and $p\gg0$, then every $q\in Q_B$ has $q\cdot a\ge
q\cdot p>0$; so $a\in D_{Q_B}$. Now let $a\notin W$. Separating $a$ from $W$ gives
$q\ne0$ with $q\cdot a\le q\cdot w$ for all $w\in W$. Then $q\ge0$, since otherwise
$w=\one+se_i$ with $q_i<0$ makes $q\cdot w\to-\infty$ as $s\to\infty$. For $b\in B_d$
and $s,t>0$, $sb+t\one\in W$; letting $t\to0$ gives $q\cdot a\le s\,q\cdot b$ for all
$s>0$, so $q\cdot b\ge0$ (let $s\to\infty$); and $t\one\in W$ gives $q\cdot a\le0$. Then
$\hat q\coloneqq q/(\one\cdot q)\in\Delta$ is non-negative on $B_d$, so $\hat q\in Q_B$,
and $\hat q\cdot a\le0$; hence $a\notin D_{Q_B}$.
\end{proof}

\begin{corollary}\label{cor:witness}
A record $(A,B)$ is safe if and only if there are $a\in A_d$, points $b_1,\dots,b_k\in
B_d$ with $0\le k\le n$, and $t_j\ge0$, such that $a-\sum_{j=1}^kt_jb_j\gg0$.
\end{corollary}

\begin{proof}
By Lemma~\ref{lem:safe}(a) it suffices to characterise $\mu(A\cap\hull(B))>0$.
$\hull(B)$ is open. If $a\in A_d$ lies in $\hull(B)$, a ball around $a$ lies in
$\hull(B)$ and meets $A$ in positive measure. Conversely, if $\mu(A\cap\hull(B))>0$,
some $a\in A_d$ lies in $\hull(B)$, since $A\setminus A_d$ is null. By
Lemma~\ref{lem:hullformula}, $a=\sum_jt_jb_j+p$ with $b_j\in B_d$, $t_j\ge0$
and $p\gg0$, and by the conic form of Carath\'eodory's theorem the sum can be rewritten
with at most $n$ linearly independent $b_j$.
\end{proof}

So a change is safe exactly when it adds some gamble $a$ that dominates a non-negative
combination of finitely many gambles that it removes.

\section{Improvement within the coherent hull}\label{sec:H}

For a model $X\in\M_U$ write
\[
  G\coloneqq\hull(X)=D_{Q_X},\qquad P\coloneqq P_{Q_X},\qquad
  Q_\tau\coloneqq(1-\tau)Q_X+\tau\Delta\quad(0<\tau<1).
\]
$Q_\tau$ is a compact convex subset of $\Delta$, and
\begin{equation}\label{eq:Ptau}
  P_{Q_\tau}(h)=(1-\tau)P(h)+\tau\min_ih_i .
\end{equation}
By Lemma~\ref{lem:hull}(b), $X\subseteq G$ a.e., so $P>0$ a.e.\ on $X$.

\subsection{The hull-preserving shrink}

\begin{lemma}\label{lem:contraction}
Let $X\in\M_U$, $c\ge0$ and $q\in\Delta$. If $q\cdot x\ge-c\,P(x)$ for $\mu$-a.e.\ $x\in
X$, then $q\in Q_X$.
\end{lemma}

\begin{proof}
For $q'\in Q_X$ put $T(q')\coloneqq(q+cq')/(1+c)\in\Delta$. For a.e.\ $x\in X$,
$(q+cq')\cdot x\ge q\cdot x+c\,P(x)\ge0$, so $T(q')\in Q_X$. Thus $T$ maps $Q_X$ into
itself, and it is a contraction with ratio $c/(1+c)<1$ and fixed point $q$. For any
$q'\in Q_X$ (which is nonempty), $T^k(q')\to q$, and $Q_X$ is closed
(Lemma~\ref{lem:hull}(a)); so $q\in Q_X$.
\end{proof}

In words: repeatedly mixing $q$ with members of the closed credal set $Q_X$ stays in
$Q_X$ and converges to $q$, so $q$ lies in $Q_X$.

\begin{proposition}\label{prop:shrink}
Let $X\in\M_U$ and $0<\tau<1$, and let $N\subseteq D_{Q_\tau}$ be measurable. Then
\[
  Q_{X\setminus N}=Q_X\qquad\text{and}\qquad\hull(X\setminus N)=\hull(X).
\]
If moreover $N\subseteq\hull(X)$ and $\mu(N\setminus X)>0$, then $N\in\M_U$ and the record
$r(N,X)$ is safe.
\end{proposition}

\begin{proof}
$Q_X\subseteq Q_{X\setminus N}$ trivially. Let $q\in Q_{X\setminus N}$ and $c\coloneqq
(1-\tau)/\tau$. If $x\in X\cap N$, then $P_{Q_\tau}(x)>0$, so by \eqref{eq:Ptau}
$\min_ix_i>-c\,P(x)$, and $q\cdot x\ge\min_ix_i>-c\,P(x)$ because $q\in\Delta$. For
a.e.\ $x\in X\setminus N$, $q\cdot x\ge0$ and $P(x)\ge0$, so $q\cdot x\ge-c\,P(x)$.
Lemma~\ref{lem:contraction} gives $q\in Q_X$.

If $N\subseteq\hull(X)$, then $N\subseteq U$ by Lemma~\ref{lem:hull}(a), and
$N\setminus X\subseteq\hull(X)=\hull(X\setminus N)$. So $\mu\bigl((N\setminus X)\cap
\hull(X\setminus N)\bigr)=\mu(N\setminus X)>0$, and $r(N,X)=(N\setminus X,X\setminus N)$
is safe by Lemma~\ref{lem:safe}(a).
\end{proof}

The identity $\hull(X\setminus N)=\hull(X)$ says that the gambles of $X$ outside $N$ still
generate, under coherence, everything that $X$ generates. For the change from $X$ to a
coherent $N\subseteq\hull(X)$ this gives more than safety: \emph{every} gamble the change
adds is implied by the gambles it removes. A coherent model that accepts what the change
removes contains $\hull(X\setminus N)\supseteq N$, so it already accepts what the change
adds, and the change cannot be applied to it.

\subsection{The repair theorem}

\begin{theorem}\label{thm:H}
Let $X\in\M_U$ be incoherent. Then there is $N\in\CC_U$ such that
\begin{enumerate}[label=(\roman*)]
  \item $m(N)-m(X)\gg0$;
  \item $N\subseteq\hull(X)$ and $\hull(X\setminus N)=\hull(X)$;
  \item $\mu(N\setminus X)>0$.
\end{enumerate}
In particular the record $r(N,X)$ is safe, and $\mu\bigl((N\setminus X)\cap
\hull(X\setminus N)\bigr)=\mu(N\setminus X)>0$.
\end{theorem}

\begin{proof}
\emph{Step 1: a strict improvement inside the hull.} Let $\M_G\coloneqq\{m(Y):Y
\subseteq G\text{ measurable}\}$. The vector measure $Y\mapsto\int_Yh\,d\mu$ on the
measurable subsets of $G$ is finite, by the first-moment condition in (M), and
non-atomic, since $\mu$ is non-atomic (Section~\ref{sec:framework}). By Lyapunov's theorem
\cite{Lyapunov,Lindenstrauss}, $\M_G$ is compact and convex. Since $X\subseteq G$ a.e.,
$x\coloneqq m(X)\in\M_G$.

For $\lambda\in\Delta$ and measurable $Y\subseteq G$,
\[
  \lambda\cdot m(Y)=\int_Y\lambda\cdot h\,d\mu\le\int_{G\cap\{\lambda\cdot h>0\}}
  \lambda\cdot h\,d\mu ,
\]
with equality if and only if $Y=G\cap\{\lambda\cdot h>0\}$ a.e., because the hyperplane
$\{\lambda\cdot h=0\}$ is null by (M). The maximiser is coherent: by
Lemma~\ref{lem:D}(a),
\[
  G\cap\{\lambda\cdot h>0\}=D_{Q_X\cup\{\lambda\}}=D_{\conv(Q_X\cup\{\lambda\})},
\]
and $\conv(Q_X\cup\{\lambda\})$ is a credal set containing $K'$.

The function $f(\lambda,z)\coloneqq\lambda\cdot(z-x)$ is bilinear and continuous on the
product of the compact convex sets $\Delta$ and $\M_G$. By the minimax theorem
\cite{Sion},
\[
  v\coloneqq\max_{z\in\M_G}\min_{\lambda\in\Delta}f(\lambda,z)
  =\min_{\lambda\in\Delta}\max_{z\in\M_G}f(\lambda,z),
\]
and both optima are attained, at $z^\ast$ and $\lambda^\ast$ say. Note that
$\min_{\lambda\in\Delta}f(\lambda,z)=\min_i(z_i-x_i)$.

Suppose $v\le0$. Then $\max_{z\in\M_G}\lambda^\ast\cdot z\le\lambda^\ast\cdot x$, so $X$
maximises $\lambda^\ast\cdot m(Y)$ over measurable $Y\subseteq G$, and $X=G\cap
\{\lambda^\ast\cdot h>0\}$ a.e. Then $X$ is coherent, a contradiction. Hence $v>0$.

Since $f(\lambda^\ast,z^\ast)\le\max_zf(\lambda^\ast,z)=v$ and $f(\lambda^\ast,z^\ast)
\ge\min_\lambda f(\lambda,z^\ast)=v$, the point $z^\ast$ maximises $\lambda^\ast\cdot z$
over $\M_G$. Writing $z^\ast=m(Y^\ast)$, uniqueness of the maximiser gives
$Y^\ast=N_0$ a.e., where
\[
  N_0\coloneqq G\cap\{\lambda^\ast\cdot h>0\}\in\CC_U,\qquad
  m(N_0)-m(X)=z^\ast-x\ge v\one\gg0 .
\]

\emph{Step 2: the gain is non-null; a uniform shrink.} Suppose $N_0\setminus X$ were
null. Then $m(N_0)-m(X)=-m(X\setminus N_0)$, and for $k\in K'$ (which is nonempty by the
standing assumption) we would have $k\cdot(m(N_0)-m(X))=-\int_{X\setminus N_0}k\cdot
h\,d\mu\le0$, since $k\cdot h>0$ on $U$. This contradicts $m(N_0)-m(X)\gg0$ and
$k\in\Delta$. So $\mu(N_0\setminus X)>0$.

For $0<\tau<1$ put $N_\tau\coloneqq N_0\cap D_{Q_\tau}$. By Lemma~\ref{lem:D}(a),
$N_\tau=D_{\conv(Q_X\cup\{\lambda^\ast\}\cup Q_\tau)}$, and this credal set contains $K'$;
so $N_\tau\in\CC_U$. Every $h\in N_0$ lies in $G$, so $P(h)>0$, and by \eqref{eq:Ptau}
$h\in D_{Q_\tau}$ for all small $\tau$. Hence $\one_{N_\tau}\to\one_{N_0}$
pointwise as $\tau\to0$, and dominated convergence, with dominating functions $|h|$ and
$1$, gives $m(N_\tau)\to m(N_0)$ and $\mu(N_0\setminus N_\tau)\to0$. Fix $\tau$ so small
that $m(N_\tau)-m(X)\gg0$ and $\mu(N_0\setminus N_\tau)<\mu(N_0\setminus X)$, and put
$N\coloneqq N_\tau$. Then (i) holds, and $\mu(N\setminus X)\ge\mu(N_0\setminus X)-
\mu(N_0\setminus N)>0$, which is (iii).

\emph{Step 3: the hull of what is removed.} Since $N\subseteq D_{Q_\tau}$,
Proposition~\ref{prop:shrink} gives $\hull(X\setminus N)=\hull(X)$, and $N\subseteq
N_0\subseteq G=\hull(X)$. This proves (ii), and the last assertions follow from
Proposition~\ref{prop:shrink}.
\end{proof}

\begin{remark}\label{rem:H}
\begin{enumerate}[label=(\alph*)]
  \item The proof uses no regularity of $X$, no pointedness of $\cl\hull(X)$, and no
  condition on the dimension of $K'$; nor does it use $K'\subseteq\relint\Delta$. It uses
  $K'\neq\emptyset$ throughout, since $Q_X\supseteq K'$ makes $\hull(X)$ and the maximisers
  coherent models on $U$ and makes Lemma~\ref{lem:contraction} applicable, and in Step 2
  through an element $k\in K'$. It uses (M) in Steps 1 and 2 and, through
  Lemma~\ref{lem:hull}(b), in Step 3.
  \item The defender is explicit: $N=\hull(X)\cap\{\lambda^\ast\cdot h>0\}\cap
  D_{(1-\tau)Q_X+\tau\Delta}$, where $\lambda^\ast$ minimises the convex function
  $\lambda\mapsto\max_{z\in\M_G}\lambda\cdot z-\lambda\cdot m(X)$ over $\Delta$, and
  $\tau>0$ is small. Every density point $a$ of $N\setminus X$ lies in $N$: a density
  point of $N\setminus X$ is a density point of $N$; a point outside $\cl N$ is not one;
  and if $a\in\partial N$, a supporting half-space of the convex set $N$ at $a$ contains
  $N$, so $N$ fills at most half of every ball around $a$. Hence $a\in N\subseteq
  \hull(X\setminus N)$, and by Lemma~\ref{lem:hullformula} and Carath\'eodory's theorem,
  as in Corollary~\ref{cor:witness}, $a-\sum_{j=1}^kt_jb_j\gg0$ for some $k\le n$,
  $b_j\in(X\setminus N)_d$ and $t_j\ge0$.
  \item The proof does not show that the record $r(N_0,X)$ produced by Step 1 is safe. The
  shrink in Step 2 is what guarantees, through Proposition~\ref{prop:shrink}, that the
  additions are implied by the removals.
  \item Step 1 applied with $G$ replaced by $U$ (which is coherent) shows that every
  $X\in\M_U$ that is not a.e.\ of the form $U\cap\{\pi\cdot h>0\}$, $\pi\in\Delta$, is
  uniformly dominated in Brier loss by a model of that form.
\end{enumerate}
\end{remark}

\section{Universal coverage}\label{sec:main}

\begin{theorem}\label{thm:main}
Let $K'\subseteq\Delta$ be nonempty, compact and convex, put $U\coloneqq D_{K'}$, and
consider the safe-improvement ranking on $\M_U$.
\begin{enumerate}[label=(\alph*)]
  \item Each $\preceq_i$ is reflexive, ties a.e.-equal models, and satisfies (G), (IB),
  (TC), (FCC) and (Add). Every nontrivial comparison is a strict Brier improvement at its
  world: if $X\preceq_iY$ and $X\ne Y$ a.e., then $m_i(X)>m_i(Y)$. There are no ties
  between models that differ on a non-null set.
  \item The undominated models are exactly the coherent models on $U$:
  $\UU=\CC^\ast=\CC_U$. No coherent model has a weak dominator other than itself.
  \item Every incoherent $X\in\M_U$ is strictly dominated by a coherent undominated model
  $N$, which moreover uniformly dominates $X$ in Brier loss, satisfies $N\subseteq
  \hull(X)=\hull(X\setminus N)$, and adds a non-null set of gambles to $X$.
  \item (GC), (ID) and (W) hold, and (C) holds in the strong form: a model tied at some
  world with a coherent model equals it.
  \item Let $K$ be a credal set. The ranking covers $K$ if and only if $K$ is
  full-dimensional and $K'\subseteq\relint K$.
\end{enumerate}
\end{theorem}

\begin{proof}
(a) is Proposition~\ref{prop:structure}.

(c) Theorem~\ref{thm:H} gives $N\in\CC_U$ with $m(N)-m(X)\gg0$, $N\subseteq\hull(X)=
\hull(X\setminus N)$, $\mu(N\setminus X)>0$, and $r(N,X)$ safe. By
Proposition~\ref{prop:dominance}, $N$ strictly dominates $X$. $N$ is undominated by
Corollary~\ref{cor:protect}.

(b) By Corollary~\ref{cor:protect}, every member of $\CC_U$ is undominated and has no
weak dominator other than itself. By (c), every incoherent model is dominated. So
$\UU=\CC_U$, and $\CC^\ast=\CC\cap\UU=\CC_U$.

(d) (GC) holds because $\UU=\CC_U$. (ID) is (c). (C): by (a), a model tied with $N$ at
some world equals $N$ a.e. (W): a coherent model weakly dominates itself, and an
incoherent model is strictly, hence weakly, dominated by a member of $\CC^\ast$, by (c).

(e) Requirements (i), (ii), (iii) and (v) of Definition~\ref{def:cover} hold by (a) and
(d). By Lemma~\ref{lem:LF}(a), $D_K\in\M_U$ if and only if $K'\subseteq K$, and then
$D_K\in\CC_U$ is undominated by (b). So the ranking covers $K$ if and only if
$K'\subseteq K$ and (LF) holds at $D_K$. If $K$ is not full-dimensional, (LF) fails by
Lemma~\ref{lem:LF}(c), since $U$ is one-sided. If $K$ is full-dimensional and
$K'\subseteq K$, (LF) holds if and only if $K'\subseteq\relint K$, by
Lemma~\ref{lem:LF}(b).
\end{proof}

By Proposition~\ref{prop:dominance}, a coherent model dominates an incoherent $X$ only if
it improves on $X$ at every world by a safe change. So, by Proposition~\ref{prop:dominance}
and Corollary~\ref{cor:protect}, (ID) for the safe-improvement ranking is equivalent to the
existence, for every incoherent $X$, of a coherent $N$ with $m(N)-m(X)\gg0$ and $r(N,X)$
safe. Theorem~\ref{thm:H} gives more, namely the hull invariance in (ii).

Since $\relint K\ne\emptyset$ exactly when $K$ is full-dimensional, the condition in (e)
can be written as $K'\subseteq\relint K$. It forces $K'\subseteq\relint\Delta$; if $K'$
meets the relative boundary of $\Delta$, no credal set is covered, although (a)--(d)
still hold.

\begin{definition}\label{def:universal}
Let $V$ be a domain. A credal set $K$ is \emph{coverable on $V$} if some ranking on $\M_V$
covers $K$. A ranking on $\M_V$ is a \emph{universal cover on $V$} if it covers every
credal set that is coverable on $V$.
\end{definition}

\begin{corollary}\label{cor:universal}
Let $K$ be a credal set. The safe-improvement ranking on $\M_U$ covers $K$ if and only if
$K$ is coverable on $U$. Hence it is a universal cover on $U$.
\end{corollary}

\begin{proof}
If some ranking on $\M_U$ covers $K$, then $D_K\in\M_U$ and (LF) holds at $D_K$. So
$K'\subseteq K$ by Lemma~\ref{lem:LF}(a), $K$ is full-dimensional by
Lemma~\ref{lem:LF}(c), and $K'\subseteq\relint K$ by Lemma~\ref{lem:LF}(b). Now apply
Theorem~\ref{thm:main}(e). The converse is trivial.
\end{proof}

\begin{corollary}\label{cor:every}
Let $K$ be a full-dimensional credal set, and let $K'\subseteq\relint K$ be nonempty,
compact and convex; for instance $K'=\{k_0\}$ with $k_0\in\relint K$. Then the
safe-improvement ranking on $\M_{D_{K'}}$ covers $K$. In particular every
full-dimensional credal set is covered by some ranking on some domain. The same ranking
covers every full-dimensional credal set whose relative interior contains $K'$.
\end{corollary}

\begin{proof}
Theorem~\ref{thm:main}(e).
\end{proof}

\subsection{Lower-dimensional credal sets}

\begin{proposition}\label{prop:lowdim}
Let $K$ be a credal set that is not full-dimensional.
\begin{enumerate}[label=(\alph*)]
  \item If $K'\subseteq K$, then $D_K$ is a coherent model on $U$, and in the
  safe-improvement ranking on $\M_U$ it is undominated and has no weak dominator other
  than itself.
  \item For every one-sided domain $V$, (LF) fails at $D_K$ for $\M_V$. Hence no ranking
  on $\M_V$ covers $K$; in particular the safe-improvement ranking on $\M_U$ does not
  cover $K$.
\end{enumerate}
\end{proposition}

\begin{proof}
(a) $D_K\in\CC_U$ by Lemma~\ref{lem:LF}(a); apply Theorem~\ref{thm:main}(b).
(b) Lemma~\ref{lem:LF}(c), and requirement (vi) of Definition~\ref{def:cover}.
\end{proof}

So the ranking does not discriminate against lower-dimensional credal sets: every
coherent model on $U$, of whatever dimension, is undominated. Lower-dimensional coherent
models on $U$ exist if and only if $K'$ is itself not full-dimensional: every credal set
containing a full-dimensional $K'$ is full-dimensional, and if $K'$ is not, then
$D_{K'}=U$ is one. For instance, if $K'=\{k_0\}$, every credal set $K$ containing $k_0$,
of any dimension, yields an undominated model $D_K$. For a lower-dimensional $K$ with
$K'\subseteq K$, requirements (i)--(v) of Definition~\ref{def:cover} hold, and exactly
requirement (vi), local freedom, fails: the boundary of $D_K$ contains a pair $\pm h$ of
nonzero gambles, and no domain that contains neighbourhoods of both can be one-sided. If
$K'\not\subseteq K$, then $D_K$ is not a model on $U$, and requirement (iv) fails as
well.

The failure of (LF) persists on every domain $V$ for which $V\cap(-V)$ contains no
non-null set of gambles with a positive and a negative coordinate
\cite[Prop.~2.10]{KCM}. Lower-dimensional targets therefore require domains that are not
one-sided, and on such domains some of them are covered by other rankings; for instance,
on the full domain the Brier ranking covers every precise model $\{p\}$ with
$p\in\relint\Delta$ \cite[Cor.~5.9]{KCM}. This note does not address whether the
safe-improvement construction yields a universal cover on domains that are not one-sided.
Definition~\ref{def:SI} makes sense on any domain, but Step 2 of the proof of
Theorem~\ref{thm:H} uses an element of $K'$, and on such domains $Q_X$ may be empty.

\section{An abstract construction lemma}\label{sec:abstract}

The argument of Sections~\ref{sec:ranking}--\ref{sec:main} separates into a construction
lemma that uses no geometry and a repair theorem (Theorem~\ref{thm:H}) that supplies its
one substantive hypothesis. The lemma is a construction theorem, not a proposed condition
on rankings: it says how to build a ranking with the dominance properties once safe repairs
are known to exist.

\smallskip\noindent\textbf{Setting.} Let $\Omega$ be a nonempty set of worlds, of any
cardinality; let $\M$ be a set of models and $\CC\subseteq\M$ a class of models called
coherent. Let $\mathcal R$ be a set, the \emph{records}, with a map $r:\M\times\M\to
\mathcal R$ and an element $0\in\mathcal R$ such that $r(X,Y)=0$ if and only if $X=Y$.
For each $\omega\in\Omega$ let $\sigma_\omega:\M\to\R$ be a score (higher is better) and
$d_\omega:\mathcal R\to\R$ a function with
\[
  d_\omega(r(X,Y))=\sigma_\omega(X)-\sigma_\omega(Y)\qquad(X,Y\in\M),
\]
and let $\mathcal G_\omega\subseteq\mathcal R$, the records that are \emph{guidance at
$\omega$}. Finally let $\mathcal T\subseteq\M^3$ be a set of triples, called
\emph{non-overlapping}. A record $\rho$ is \emph{safe} if $\rho\ne r(Y,C)$ for all
$Y\in\M$ and $C\in\CC$. Define
\[
  X\preceq_\omega Y\iff r(X,Y)\in\mathcal G_\omega,\ \text{or $r(X,Y)$ is safe and }
  d_\omega(r(X,Y))>0 .
\]
Consider the hypotheses:
\begin{itemize}
  \item[(A1)] For every $\omega$, $0\in\mathcal G_\omega$, and $d_\omega(\rho)>0$ for every
  $\rho\in\mathcal G_\omega\setminus\{0\}$.
  \item[(A2)] If $C\in\CC$, $Y\in\M$ and $r(Y,C)\in\mathcal G_\omega$ for every $\omega$,
  then $Y=C$.
  \item[(A3)] For every $X\in\M\setminus\CC$ there is $N\in\CC$ such that $r(N,X)$ is safe
  and $\sigma_\omega(N)>\sigma_\omega(X)$ for every $\omega$.
  \item[(A4)] For every $(X,Y,Z)\in\mathcal T$ and every $\omega$: if $r(X,Y),r(Y,Z)\in
  \mathcal G_\omega$, then $r(X,Z)\in\mathcal G_\omega$; and if $r(X,Y)$ or $r(Y,Z)$ is
  safe, then $r(X,Z)$ is safe.
  \item[(A5)] Every record $\rho\ne0$ with $\rho\in\mathcal G_\omega$ for every $\omega$
  is safe.
\end{itemize}
(A1) says that guidance strictly improves the score of its own world; (A2) that guidance
alone cannot weakly dominate a coherent model nontrivially; (A3) that every incoherent
model has a coherent repair that improves its score at every world by a safe change; and
(A4) that guidance composes over non-overlapping steps and that safety cannot be lost by
adjoining a further non-overlapping step; (A5) is used only for a converse. Safety is
defined here through the record map: a record is unsafe exactly when it describes a change
out of some coherent model. Upward closure of safety is not derived from this definition
but assumed in (A4); in the application it holds because every record that applies to a
model is the record of a change out of it (Lemma~\ref{lem:record}(a)).

\begin{lemma}\label{lem:abstract}
In the setting above, with weak dominance, dominance and strict dominance defined as in
Definition~\ref{def:dom}:
\begin{enumerate}[label=(\alph*)]
  \item Assume (A1). Each $\preceq_\omega$ is reflexive and depends on $(X,Y)$ only through
  $r(X,Y)$. If $X\preceq_\omega Y$ and $X\ne Y$, then $\sigma_\omega(X)>\sigma_\omega(Y)$;
  hence $X$ and $Y$ are tied at $\omega$ only if $X=Y$, and $\preceq_\omega$ satisfies
  (FCC) and (TC). If $r(X,Y)\in\mathcal G_\omega\setminus\{0\}$, then $X\prec_\omega Y$.
  \item Assume (A1) and (A2). No $C\in\CC$ has a weak dominator other than itself; in
  particular every coherent model is undominated.
  \item Assume (A1) and (A3). Every $X\in\M\setminus\CC$ is strictly dominated by some
  $N\in\CC$.
  \item Assume (A1)--(A3). The undominated models are exactly the coherent models; every
  model is weakly dominated by an undominated coherent model; and a model tied at some
  world with a coherent model equals it.
  \item Assume (A1) and (A4). If $(X,Y,Z)\in\mathcal T$, $X\preceq_\omega Y$ and
  $Y\preceq_\omega Z$, then $X\preceq_\omega Z$.
  \item Assume (A1) and (A5). If $N$ weakly dominates $X$ and $N\ne X$, then $r(N,X)$ is
  safe and $\sigma_\omega(N)>\sigma_\omega(X)$ for every $\omega$. Hence, under (A1) and
  (A5), (A3) holds if and only if every model outside $\CC$ is dominated by a member of
  $\CC$.
\end{enumerate}
\end{lemma}

\begin{proof}
(a) Reflexivity holds because $r(X,X)=0\in\mathcal G_\omega$. If $X\preceq_\omega Y$ and
$X\ne Y$, then $\rho\coloneqq r(X,Y)\ne0$ is accepted either as guidance, and then
$d_\omega(\rho)>0$ by (A1), or as a safe record with $d_\omega(\rho)>0$; in both cases
$\sigma_\omega(X)-\sigma_\omega(Y)=d_\omega(\rho)>0$. A tie between $X\ne Y$ would give
$\sigma_\omega(X)>\sigma_\omega(Y)>\sigma_\omega(X)$. A chain $A_0\preceq_\omega\dots
\preceq_\omega A_k$ gives $\sigma_\omega(A_0)\ge\sigma_\omega(A_k)$, while $A_k
\prec_\omega A_0$ would give $A_k\ne A_0$ (by reflexivity) and so $\sigma_\omega(A_k)>
\sigma_\omega(A_0)$; this proves (FCC), and (TC) is its case $k=2$. If $r(X,Y)\in
\mathcal G_\omega\setminus\{0\}$, then $X\preceq_\omega Y$ and $X\ne Y$, so they are not
tied and $X\prec_\omega Y$.

(b) Let $Y$ weakly dominate $C\in\CC$, and put $\rho\coloneqq r(Y,C)$. By definition
$\rho$ is not safe, so $Y\preceq_\omega C$ requires $\rho\in\mathcal G_\omega$, for every
$\omega$. By (A2), $Y=C$.

(c) Take $N$ as in (A3). For every $\omega$, $r(N,X)$ is safe and $d_\omega(r(N,X))=
\sigma_\omega(N)-\sigma_\omega(X)>0$, so $N\preceq_\omega X$; and $N\ne X$ because
$N\in\CC$ and $X\notin\CC$. By (a), $N\prec_\omega X$.

(d) By (b) and (c). For the last claim use (a).

(e) If $r(X,Y),r(Y,Z)\in\mathcal G_\omega$, then $r(X,Z)\in\mathcal G_\omega$ by (A4).
Otherwise one of $r(X,Y)$, $r(Y,Z)$ is safe with positive $d_\omega$; then $r(X,Z)$ is
safe by (A4), and $d_\omega(r(X,Z))=\bigl(\sigma_\omega(X)-\sigma_\omega(Y)\bigr)+
\bigl(\sigma_\omega(Y)-\sigma_\omega(Z)\bigr)>0$, because one summand is positive and the
other is non-negative by (a).

(f) $\rho\coloneqq r(N,X)\ne0$ is accepted at every $\omega$, so $\sigma_\omega(N)>
\sigma_\omega(X)$ for every $\omega$ by (a). If $\rho$ were not safe, it would be accepted
as guidance at every $\omega$, and (A5) would make it safe. For the equivalence: (c) gives
one direction; conversely, if $N\in\CC$ dominates $X\notin\CC$, then $N\ne X$ and $N$
weakly dominates $X$, so $N$ is as required in (A3).
\end{proof}

The record layer is what gives the lemma content. Without records one could let
$X\preceq_\omega Y$ hold when $r(X,Y)\in\mathcal G_\omega$, or when $Y\notin\CC$ and
$\sigma_\omega(X)>\sigma_\omega(Y)$. Given coherent models that improve every incoherent
model at every world, this yields protection and dominance. But acceptance then depends on
whether $Y$ is coherent, which is not a property of $r(X,Y)$. In the setting of
Example~\ref{ex:tr}, with $C=U$ and $A,B$ as there, the model $C\setminus A$ is incoherent:
if it were a.e.\ equal to some $D_Q$, then $D_Q$ would contain the open set $U\setminus\cl
A$ (Lemma~\ref{lem:D}(d)), whose convex hull contains $A$, while $\mu(A\cap D_Q)=0$. So, at
world 1, the record $(\emptyset,B)$ is refused from $C$ (it is not guidance there, and $C$
is coherent) but accepted from $C\setminus A$; hence (IB) fails. Moreover $C\setminus(A\cup
B)\preceq_1C\setminus A\preceq_1C$ by steps with the disjoint supports $B$ and $A$, while
$C\setminus(A\cup B)\preceq_1C$ fails; hence (Add) fails. Safety is a record-level
substitute for the condition that the change starts from an incoherent model.

The hypotheses involve no topology, measure or finiteness. In the setting of
Sections~\ref{sec:framework}--\ref{sec:main}, take $\M=\M_U$, $\CC=\CC_U$, the records of
Section~\ref{sec:ranking} with $0$ the null record, $\sigma_i=m_i$, $d_i(r)=m_i(r)$,
$\mathcal G_i=G_i$, and $\mathcal T$ the triples with $(X\triangle Y)\cap(Y\triangle Z)$
null. By Lemma~\ref{lem:record}(a), safety in the abstract sense is safety in the sense of
Section~\ref{sec:ranking}, and $\preceq_\omega$ is the safe-improvement ranking. (A1) is
Lemma~\ref{lem:increment}(a). (A2) holds because a record $r(Y,C)$ in every $G_i$ has the
form $(A,\emptyset)$ with $A\subseteq\Rpp\subseteq C$ (Lemma~\ref{lem:increment}(c)),
while $A=Y\setminus C$ misses $C$; so $A$ is null and $Y=C$ a.e. (A3) is
Theorem~\ref{thm:H}. (A4) follows from Lemma~\ref{lem:record}(b), the closure of $G_i$
under $\sqcup$, and Lemma~\ref{lem:safe}(c). (A5) is Lemma~\ref{lem:increment}(c).
Lemma~\ref{lem:abstract} thus yields the order-theoretic and dominance clauses of
Theorem~\ref{thm:main}(a)--(d), apart from (IB) and the form of (G) used in
Definition~\ref{def:cover}, which are Proposition~\ref{prop:structure}(a),(c). The
geometric clauses of Theorem~\ref{thm:main}(c) come from Theorem~\ref{thm:H}, and uniform
dominance in Brier loss uses the finiteness of $\Omega$.

\section{Infinitely many worlds}\label{sec:infinite}

This section extends Theorem~\ref{thm:H}, the safe-improvement ranking and the coverage
results to arbitrary sets of worlds. Lyapunov's theorem is not available in infinite
dimensions; it is replaced by a relaxation of acceptance sets that is used only inside
the proof. The finite setting of Sections~\ref{sec:framework}--\ref{sec:main} is the
special case $\Omega=\{1,\dots,n\}$, $E=\R^n$ (Remark~\ref{rem:finite-special}).

\subsection{Gambles and states}

Let $\Omega$ be a nonempty set, of any cardinality, and let $\ell^\infty(\Omega)$ be the
space of bounded real functions on $\Omega$ with the norm $\|h\|\coloneqq\sup_\Omega|h|$
and the pointwise order. Throughout this section $E\subseteq\ell^\infty(\Omega)$ is a
norm-closed, \emph{separable} linear subspace with $\one\in E$; its elements are the
\emph{gambles}. (Remark~\ref{rem:orderunit} replaces $\one$ by any order unit bounded
below by a positive constant.) A \emph{state} is a linear functional $q:E\to\R$ that is
positive ($q(h)\ge0$ whenever $h\ge0$) and normalised ($q(\one)=1$); $S$ is the set of
states, with the weak-$\ast$ topology $\sigma(E^\ast,E)$. For $\omega\in\Omega$,
$\delta_\omega(h)\coloneqq h(\omega)$ defines a state. A \emph{credal set} is a nonempty
weak-$\ast$ closed convex subset of $S$. For a nonempty weak-$\ast$ compact $Q\subseteq S$
put
\[
  P_Q(h)\coloneqq\min_{q\in Q}q(h),\qquad D_Q\coloneqq\{h\in E:P_Q(h)>0\},\qquad
  C_Q\coloneqq\{h\in E:P_Q(h)\ge0\};
\]
the minimum is attained because $q\mapsto q(h)$ is weak-$\ast$ continuous.

\begin{lemma}\label{lem:states}
\begin{enumerate}[label=(\alph*)]
  \item Every state has norm $1$. $S$ is convex, weak-$\ast$ compact and metrizable.
  \item $P_S(h)=\inf_{\omega\in\Omega}h(\omega)$ for every $h\in E$.
  \item Let $\hat\Omega$ be the weak-$\ast$ closure of $\{\delta_\omega:\omega\in\Omega\}$
  in $S$, a compact metrizable space. The map $h\mapsto\hat h$, $\hat h(s)\coloneqq s(h)$,
  is a linear isometry of $E$ into $C(\hat\Omega)$ that preserves and reflects the order
  and maps $\one$ to $1$. Every state $q$ is of the form $q(h)=\int_{\hat\Omega}\hat
  h\,d\pi_q$ for some regular Borel probability measure $\pi_q$ on $\hat\Omega$.
\end{enumerate}
\end{lemma}

\begin{proof}
(a) For $q\in S$ and $h\in E$, $-\|h\|\one\le h\le\|h\|\one$ gives $|q(h)|\le\|h\|$, with
equality at $h=\one$. So $S$ lies in the unit ball of $E^\ast$, and it is weak-$\ast$
closed, since positivity and normalisation are pointwise conditions. By the
Banach--Alaoglu theorem it is weak-$\ast$ compact, and it is metrizable because bounded
subsets of the dual of a separable space are weak-$\ast$ metrizable. Convexity is clear.

(b) $\delta_\omega\in S$ gives $P_S(h)\le\inf_\Omega h$. Conversely $h-(\inf_\Omega
h)\one\ge0$ lies in $E$, so $q(h)\ge\inf_\Omega h$ for every $q\in S$.

(c) $\hat\Omega$ is a closed subset of the compact metrizable space $S$. Each $\hat h$ is
continuous on $\hat\Omega$. By (a), $|\hat h|\le\|h\|$ on $S$, and $\hat h(\delta_\omega)=
h(\omega)$; so $\sup_{\hat\Omega}|\hat h|=\|h\|$. If $h\ge0$, then $\hat h\ge0$ on $S$;
if $\hat h\ge0$ on $\hat\Omega$, then $h(\omega)=\hat h(\delta_\omega)\ge0$. Let $q\in S$.
On the subspace $\hat E\coloneqq\{\hat h:h\in E\}$ of $C(\hat\Omega)$ put $L(\hat h)
\coloneqq q(h)$, which is well defined because $h\mapsto\hat h$ is injective. Since
$(\sup_\Omega h)\one-h\ge0$, $L(\hat h)\le\sup_\Omega h=\max_{\hat\Omega}\hat h$. The
Hahn--Banach theorem extends $L$ to a linear functional on $C(\hat\Omega)$ dominated by
the sublinear functional $f\mapsto\max_{\hat\Omega}f$. The extension is positive, since
$f\ge0$ gives $-L(f)=L(-f)\le\max(-f)\le0$, and $L(1)=q(\one)=1$ because $1=\hat\one\in
\hat E$. The Riesz representation theorem gives $\pi_q$.
\end{proof}

So, whatever $\Omega$ is, the states are countably additive probabilities on the
compact metrizable space $\hat\Omega$, though in general not on $\Omega$, and the
representing measure need not be unique. If $\Omega$ is a compact metric space and $E=C(\Omega)$,
then the states are exactly the regular Borel probability measures on $\Omega$, by the Riesz
representation theorem.

\begin{lemma}\label{lem:envelopes}
Let $Q,Q_1,Q_2$ be credal sets.
\begin{enumerate}[label=(\alph*)]
  \item $P_Q$ is $1$-Lipschitz, concave and positively homogeneous, so $D_Q$ is a
  norm-open convex cone and $C_Q$ is norm-closed. $D_S\subseteq D_Q$, $D_Q\cap(-D_Q)=
  \emptyset$, and $D_{Q_1}\cap D_{Q_2}=D_{\conv(Q_1\cup Q_2)}$, where $\conv(Q_1\cup Q_2)$
  is a credal set.
  \item $P_Q(h+t\one)=P_Q(h)+t$. The norm boundary $\partial D_Q$ is contained in
  $\{P_Q=0\}$, and $\cl D_Q\subseteq C_Q$.
  \item For $0\le\tau\le1$, $(1-\tau)Q+\tau S$ is a credal set, and
  $P_{(1-\tau)Q+\tau S}=(1-\tau)P_Q+\tau P_S$.
  \item If $D_{Q_1}\subseteq D_{Q_2}$, then $Q_2\subseteq Q_1$.
\end{enumerate}
\end{lemma}

\begin{proof}
(a) Each $q\in Q$ has norm $1$, so $|P_Q(h)-P_Q(h')|\le\|h-h'\|$; a minimum of linear
functions is concave and positively homogeneous. If $h\in D_S$ then $q(h)\ge P_S(h)>0$
for all $q\in Q\subseteq S$. A linear functional has the same minimum over $Q_1\cup Q_2$
and over its convex hull, which is weak-$\ast$ compact as a continuous image of
$Q_1\times Q_2\times[0,1]$.
(b) The identity holds because $q(\one)=1$. If $P_Q(h)<0$ or $P_Q(h)>0$, a neighbourhood
of $h$ lies outside $D_Q$ or inside it, by (a). $C_Q$ is closed and contains $D_Q$.
(c) The set is a continuous image of $Q\times S$, hence compact, and it is convex. The
minimum of $(1-\tau)q+\tau s$ over the product is the sum of the minima.
(d) Suppose $k\in Q_2\setminus Q_1$. The continuous linear functionals on $(E^\ast,
\sigma(E^\ast,E))$ are the evaluations at elements of $E$, so strict separation of $k$
from the compact convex set $Q_1$ gives $g\in E$ and $c\in\R$ with $q(g)>c>k(g)$ for all
$q\in Q_1$. Put $h\coloneqq g-c\one$. Then $P_{Q_1}(h)>0$, because the minimum over $Q_1$
is attained, while $k(h)<0$. So $h\in D_{Q_1}\setminus D_{Q_2}$.
\end{proof}

\subsection{The measure}

Let $\mu$ be a finite Borel measure on $(E,\|\cdot\|)$. We use three conditions.
\begin{itemize}
  \item[(M1)] $\int\|h\|\,d\mu(h)<\infty$.
  \item[(M2)] $\mu\{P_Q=0\}=0$ for every credal set $Q$.
  \item[(M3)] $\mu$ has full support: every nonempty norm-open set has positive measure.
\end{itemize}
(M2) applies in particular to $Q=\{q\}$ for a state $q$, where $\{P_Q=0\}=\{q(h)=0\}$; to
$Q=\{\delta_\omega\}$, where it says that $\{h(\omega)=0\}$ is null; and to $Q=S$, where
it says that the gambles $h$ with $\inf_\Omega h=0$ form a null set. Equalities and
inclusions between sets of gambles hold a.e.\ if they hold up to a $\mu$-null set.

\begin{lemma}\label{lem:measure}
Let $\omega_0\in\Omega$ and $E_0\coloneqq\{g\in E:g(\omega_0)=0\}$. Let $\nu$ be a Borel
probability measure on $E_0$ with full support and $\int\|g\|\,d\nu<\infty$; for instance
$\nu=\sum_kw_k\delta_{g_k}$, where $(g_k)$ is a dense sequence in $E_0$ and $w_k>0$,
$\sum_kw_k=1$, $\sum_kw_k(1+\|g_k\|)<\infty$. Let $\lambda$ be a Borel probability measure
on $\R$ with full support, no atoms and $\int|t|\,d\lambda<\infty$; for instance the
standard normal distribution. Let $\mu$ be the image of $\nu\otimes\lambda$ under
$\Phi(g,t)\coloneqq g+t\one$. Then $\mu$ is a Borel probability measure on $E$ satisfying
(M1), (M2) and (M3). Moreover $\mu\{P_Q=0\}=0$ for every nonempty weak-$\ast$ compact
$Q\subseteq S$.
\end{lemma}

\begin{proof}
$E_0$ is a closed subspace of the separable space $E$, so it is separable, and dense
sequences exist. $\Phi:E_0\times\R\to E$ is linear and continuous, and it is bijective
with continuous inverse $h\mapsto(h-h(\omega_0)\one,h(\omega_0))$; so it is a linear
homeomorphism. Since $E_0$ and $\R$ are separable metric spaces, the Borel
$\sigma$-algebra of $E_0\times\R$ is the product $\sigma$-algebra, so $\mu$ is a Borel
probability measure on $E$.

(M1): $\int\|h\|\,d\mu=\iint\|g+t\one\|\,d\nu\,d\lambda\le\int\|g\|\,d\nu+\int|t|\,
d\lambda<\infty$. For the atomic example, $\int\|g\|\,d\nu=\sum_kw_k\|g_k\|<\infty$, and
$\nu$ has full support because every nonempty open subset of $E_0$ contains some $g_k$.

(M3): if $W\subseteq E$ is nonempty and open, then $\Phi^{-1}(W)$ is nonempty and open,
so it contains a product $V\times I$ of nonempty open sets, and $\mu(W)\ge\nu(V)\lambda(I)
>0$.

(M2): let $Q\subseteq S$ be nonempty and weak-$\ast$ compact. $P_Q$ is continuous, and
$P_Q(g+t\one)=P_Q(g)+t$. So $\Phi^{-1}\{P_Q=0\}=\{(g,t):t=-P_Q(g)\}$, and by Fubini's
theorem $\mu\{P_Q=0\}=\int\lambda(\{-P_Q(g)\})\,d\nu(g)=0$, since $\lambda$ has no atoms.
\end{proof}

With the atomic $\nu$, $\mu$ is carried by the countably many lines $g_k+\R\one$. This is
admissible; but it shows that the null sets of $\mu$ are a substantive choice
(Section~\ref{sec:scope}).

\smallskip\noindent\textbf{Where the conditions are used.} (M1) is used for the Bochner
integrals below and for dominated convergence. (M2) is used for the uniqueness of
maximisers and for the hull lemma (with $Q=\{q\}$ and $Q=Q_B$), for strict guidance (with
$Q=\{\delta_\omega\}$) and for protection (with $Q=S$). (M3) is used only to identify
models: to show that the credal set of a coherent model is unique, that a coherent model
lies in $U$ exactly when its credal set contains $K'$, that $U$ is one-sided in the sense
of Section~\ref{subsec:inf-LF}, and to characterise local freedom.
The existence of the defender in Theorem~\ref{thm:Hinf} and the protection of coherent
models do not use (M3).

\smallskip
The condition $\mu\{P_S=0\}=0$ in (M2) cannot be dropped: without it, guidance alone can
strictly dominate a coherent model.

\begin{example}\label{ex:PS}
Let $\Omega=(0,1]$ and $E\coloneqq\{a+b\,\omega:a,b\in\R\}$, a two-dimensional subspace of
$\ell^\infty(\Omega)$ containing $\one$. Every state has the form $q_s(a+b\,\omega)=a+bs$
with $s\in[0,1]$ (positivity applied to $\omega$ and $\one-\omega$), and $P_S(a+b\,\omega)
=\min(a,a+b)$. The state $q_0$ is not an evaluation. Let $K'\coloneqq\{q_{1/2}\}$, so that
$U=\{a+b/2>0\}$, and $Q\coloneqq\{q_s:0\le s\le\tfrac12\}$, so that $C\coloneqq D_Q=\{a>0,
\ a+b/2>0\}$ is a coherent model on $U$. Let $\mu$ be the sum of a standard Gaussian
distribution on the coefficients $(a,b)$ and the distribution of $(0,\beta)$ with $\beta$
exponentially distributed. Then $\mu$ is finite, has full support and satisfies (M1), but
the ray $R\coloneqq\{a=0,\ b>0\}$ has $\mu(R)=1$ and lies in $\{P_S=0\}$. $R\subseteq U$
and $R\cap C=\emptyset$. Every $h=b\,\omega\in R$ is positive at every world, so the record
$(R,\emptyset)$ is guidance at every $\omega\in\Omega$ and non-null, and $m_\omega(R)=
\int_Rh(\omega)\,d\mu>0$. With the definitions of Section~\ref{subsec:inf-ranking},
$C\cup R\preceq_\omega C$ for every $\omega$, while $C\preceq_\omega C\cup R$ fails, because
the record $(\emptyset,R)$ is not guidance at $\omega$ and has negative increment. So the
model $C\cup R$ strictly dominates the coherent model $C$ at every world. Note that
$C\cup R$ is a convex cone that excludes $0$ and contains every $h$ with
$\inf_\Omega h>0$.
\end{example}

\subsection{Models and coherence}

Fix a credal set $K'$ and put $U\coloneqq D_{K'}$, a norm-open convex cone with
$U\cap(-U)=\emptyset$, containing no $h$ with $\sup_\Omega h\le0$. Models on $U$ are the
Borel sets $X\subseteq U$, identified when they agree a.e.; their class is $\M_U$. All the
sets $D_Q$ are norm-open by Lemma~\ref{lem:envelopes}(a), hence Borel; so are the sets
$\hull(X)$, $N_0$ and $N_\tau$ below. A model $X$ is \emph{coherent} if $X=D_Q$ a.e.\ for
a credal set $Q$. Put
\[
  \CC_U\coloneqq\{D_Q:Q\text{ a credal set with }K'\subseteq Q\}\subseteq\M_U .
\]
For a Borel set $B\subseteq U$ let
\[
  Q_B\coloneqq\{q\in S:q(h)\ge0\text{ for $\mu$-a.e.\ }h\in B\},\qquad\hull(B)\coloneqq
  D_{Q_B}.
\]
If $B$ is null, then $Q_B=S$ and $\hull(B)=D_S$.

\begin{lemma}\label{lem:inf-hull}
Assume (M2), and let $B\subseteq U$ be Borel.
\begin{enumerate}[label=(\alph*)]
  \item $Q_B$ is a credal set containing $K'$; so $\hull(B)\in\CC_U$.
  \item $P_{Q_B}\ge0$ a.e.\ on $B$, and $B\subseteq\hull(B)$ a.e.
  \item For every credal set $Q$ the following are equivalent: $B\subseteq D_Q$ a.e.;
  $Q\subseteq Q_B$; $\hull(B)\subseteq D_Q$. So $\hull(B)$ is the smallest set of the form
  $D_Q$ containing $B$ a.e.
  \item Assume also (M3). If $O$ is norm-open and $O\subseteq D_Q$ a.e., then $O\subseteq
  D_Q$. Consequently $D_Q\subseteq U$ a.e.\ if and only if $K'\subseteq Q$; $D_{Q_1}=
  D_{Q_2}$ a.e.\ only if $Q_1=Q_2$; and $\CC_U$ is the class of coherent models on $U$,
  each with a unique credal set.
\end{enumerate}
\end{lemma}

\begin{proof}
(a) $K'\subseteq Q_B$ because $k(h)>0$ on $U$ for $k\in K'$, and $Q_B$ is convex. Since
$S$ is metrizable, closedness can be tested along sequences. Let $q_j\in Q_B$ converge
weak-$\ast$ to $q$. If $q(h)<0$, then $q_j(h)<0$ for some $j$; so $B\cap\{q<0\}\subseteq
\bigcup_jB\cap\{q_j<0\}$, which is null, and $q\in Q_B$.
(b) $Q_B$ is a compact metrizable space, so it has a dense sequence $(q_j)$. Off a null
set, every $h\in B$ satisfies $q_j(h)\ge0$ for all $j$, and then $q(h)\ge0$ for all
$q\in Q_B$, by weak-$\ast$ continuity of $q\mapsto q(h)$; that is, $P_{Q_B}(h)\ge0$. The
set $\{P_{Q_B}=0\}$ is null by (M2), so $B\subseteq D_{Q_B}$ a.e.
(c) If $B\subseteq D_Q$ a.e., each $q\in Q$ is positive a.e.\ on $B$, so $Q\subseteq Q_B$.
If $Q\subseteq Q_B$, then $P_{Q_B}\le P_Q$ and $D_{Q_B}\subseteq D_Q$. If $\hull(B)
\subseteq D_Q$, then $B\subseteq D_Q$ a.e.\ by (b).
(d) The set $O\cap\{P_Q<0\}$ is open and null, hence empty by (M3). If $h\in O$, then
$h-\eta\one\in O$ for some $\eta>0$, so $P_Q(h)-\eta\ge0$ and $h\in D_Q$. If $D_Q
\subseteq U$ a.e., then $D_Q\subseteq U=D_{K'}$, and Lemma~\ref{lem:envelopes}(d) gives
$K'\subseteq Q$; conversely, $K'\subseteq Q$ gives $P_Q\le P_{K'}$, so $D_Q\subseteq U$. If $D_{Q_1}=D_{Q_2}$
a.e., the first claim gives $D_{Q_1}=D_{Q_2}$, and Lemma~\ref{lem:envelopes}(d) gives
$Q_1=Q_2$.
\end{proof}

Records, their application to models, the operation $\sqcup$ and safety are defined as in
Section~\ref{sec:ranking}, with Borel subsets of $U$ in place of measurable subsets of
$U\subseteq\R^n$: a record is \emph{safe} if it applies to no member of $\CC_U$.
Lemma~\ref{lem:record} and Lemma~\ref{lem:safe}(c) are set-theoretic and hold verbatim.

\begin{lemma}\label{lem:inf-safe}
Assume (M2), and let $r=(A,B)$ be a record. The following are equivalent: (i) $r$ is
safe; (ii) $r$ applies to no set $D_Q$ with $Q$ a credal set; (iii)
$\mu(A\cap\hull(B))>0$. If $B$ is null, $r$ is safe if and only if $\mu(A\cap D_S)>0$.
\end{lemma}

\begin{proof}
(ii)$\Rightarrow$(i) is trivial. (i)$\Rightarrow$(iii): if $A\cap\hull(B)$ is null, then
$r$ applies to $\hull(B)\in\CC_U$, by Lemma~\ref{lem:inf-hull}(a),(b).
(iii)$\Rightarrow$(ii): if $r$ applies to $D_Q$, then $B\subseteq D_Q$ a.e., so
$\hull(B)\subseteq D_Q$ by Lemma~\ref{lem:inf-hull}(c), and $A\cap\hull(B)$ is null. The
last claim holds because $\hull(B)=D_S$ when $B$ is null.
\end{proof}

\subsection{The repair theorem}

By (M1), and because $E$ is separable, the identity map of $E$ is Bochner integrable with
respect to $\mu$. For a Borel set $Y\subseteq E$ put $m(Y)\coloneqq\int_Yh\,d\mu(h)\in E$.
Bounded linear functionals commute with the Bochner integral, so $q(m(Y))=\int_Yq(h)\,
d\mu$ for $q\in S$, and in particular $m(Y)(\omega)=\int_Yh(\omega)\,d\mu$; also $\|m(Y)
\|\le\int_Y\|h\|\,d\mu$.

\begin{theorem}\label{thm:Hinf}
Assume (M1) and (M2). Let $X\in\M_U$ with $X\notin\CC_U$. Then there is $N\in\CC_U$ such
that
\begin{enumerate}[label=(\roman*)]
  \item $\inf_{q\in S}q\bigl(m(N)-m(X)\bigr)>0$; in particular $\inf_{\omega\in\Omega}
  \bigl(m(N)(\omega)-m(X)(\omega)\bigr)>0$;
  \item $N\subseteq\hull(X)$ and $\hull(X\setminus N)=\hull(X)$;
  \item $\mu(N\setminus X)>0$.
\end{enumerate}
In particular the record $r(N,X)$ is safe.
\end{theorem}

The margin in (i) holds at every state, not merely at every world; by
Lemma~\ref{lem:states}(b) the two infima in (i) are equal.

\begin{proof}
\emph{Step A: the hull.} Put $G\coloneqq\hull(X)=D_{Q_X}$ and $P\coloneqq P_{Q_X}$. By
Lemma~\ref{lem:inf-hull}(a),(b), $Q_X$ is a credal set containing $K'$, $G\in\CC_U$, $P\ge0$
a.e.\ on $X$, and $X\subseteq G$ a.e.

\emph{Step B: relaxed choices and the minimax theorem.} Since $\mu$ is finite,
$L^\infty(\mu)$ is the dual of $L^1(\mu)$. The set $F_G\coloneqq\{f\in L^\infty(\mu):0\le
f\le\one_G\text{ a.e.}\}$ is convex, bounded and weak-$\ast$ closed (it is cut out by the
conditions $\int f\varphi\,d\mu\ge0$ and $\int(\one_G-f)\varphi\,d\mu\ge0$ for $0\le
\varphi\in L^1(\mu)$), hence weak-$\ast$ compact. Define $T:F_G\to E$ by $T(f)\coloneqq\int
f(h)\,h\,d\mu(h)$, so that $T(\one_Y)=m(Y)$ for Borel $Y\subseteq G$. $T$ is
weak-$\ast$-to-norm continuous: the identity map lies in $L^1(\mu;E)$, so there are simple
functions $s_n=\sum_ie_{n,i}\one_{E_{n,i}}$ with $\int\|h-s_n(h)\|\,d\mu<1/n$; the maps
$T_n(f)\coloneqq\sum_ie_{n,i}\int_{E_{n,i}}f\,d\mu$ are weak-$\ast$-to-norm continuous,
and $\|T(f)-T_n(f)\|\le\|f\|_\infty/n\le1/n$ on $F_G$, so $T$ is a uniform limit of
continuous maps. Hence $Z_G\coloneqq T(F_G)$ is norm-compact and convex. By Step A,
$\one_X\in F_G$, so $x\coloneqq m(X)\in Z_G$. (We do not claim that $\{m(Y):Y\subseteq G\}$
is convex; relaxation is a device of the proof, and the defender will be an ordinary
model.)

For $q\in S$ and $f\in F_G$, $q(T(f))=\int f(h)q(h)\,d\mu\le\int_{G\cap\{q>0\}}q(h)\,d\mu$,
with equality if and only if $f=1$ a.e.\ on $G\cap\{q>0\}$ and $f=0$ a.e.\ on $G\cap\{q<0\}$,
that is, since $\{q=0\}$ is null by (M2), if and only if $f=\one_{G\cap\{q>0\}}$ a.e. The
maximiser is a member of $\CC_U$: $G\cap\{q>0\}=D_{\conv(Q_X\cup\{q\})}$ by
Lemma~\ref{lem:envelopes}(a), and this credal set contains $K'$.

Consider $\Phi(q,z)\coloneqq q(z-x)$ on $S\times Z_G$. $S$ is a compact convex subset of
the locally convex space $(E^\ast,\sigma(E^\ast,E))$ and $Z_G$ a compact convex subset of
$(E,\|\cdot\|)$. $\Phi$ is affine in each variable, weak-$\ast$ continuous in $q$ for fixed
$z$, and norm-continuous in $z$ for fixed $q$. By Sion's minimax theorem \cite{Sion},
\[
  v\coloneqq\sup_{z\in Z_G}\min_{q\in S}\Phi(q,z)=\min_{q\in S}\sup_{z\in Z_G}\Phi(q,z).
\]
Both optima are attained: $z\mapsto\min_S\Phi(q,z)=\inf_\Omega(z-x)$ is $1$-Lipschitz on
the compact set $Z_G$, so it attains its maximum at some $z^\ast$; and $q\mapsto\max_{Z_G}
\Phi(q,z)$ is weak-$\ast$ lower semicontinuous, as a supremum of continuous functions, so
it attains its minimum on $S$ at some $q^\ast$.

Since $x\in Z_G$, $v\ge\min_S\Phi(q,x)=0$. If $v=0$, then $q^\ast(z)\le q^\ast(x)$ for
all $z\in Z_G$, so $\one_X$ maximises $f\mapsto q^\ast(T(f))$ over $F_G$, and $X=G\cap
\{q^\ast>0\}$ a.e.; then $X\in\CC_U$, a contradiction. Hence $v>0$.

Moreover $v=\min_S\Phi(\cdot,z^\ast)\le\Phi(q^\ast,z^\ast)\le\max_{Z_G}\Phi(q^\ast,\cdot)=
v$, so $z^\ast$ maximises $q^\ast$ over $Z_G$. Write $z^\ast=T(f^\ast)$. Then $f^\ast$
maximises $f\mapsto q^\ast(T(f))$, so $f^\ast=\one_{N_0}$ a.e., where
\[
  N_0\coloneqq G\cap\{q^\ast>0\}\in\CC_U,
\]
and $z^\ast=m(N_0)$. Consequently $q(m(N_0)-m(X))\ge\min_S\Phi(\cdot,z^\ast)=v$ for every
$q\in S$, and in particular $m(N_0)(\omega)-m(X)(\omega)\ge v$ for every $\omega\in\Omega$.

\emph{Step C: the shrink.} For $0<\tau<1$ let $Q_\tau\coloneqq(1-\tau)Q_X+\tau S$, a
credal set with $P_{Q_\tau}=(1-\tau)P+\tau P_S$ (Lemma~\ref{lem:envelopes}(c)), and put
$N_\tau\coloneqq N_0\cap D_{Q_\tau}=D_{\conv(Q_X\cup\{q^\ast\}\cup Q_\tau)}\in\CC_U$.
Since $P\ge P_S$, $P_{Q_\tau}=P-\tau(P-P_S)$ does not increase with $\tau$; and each
$h\in N_0$ has $P(h)>0$, hence lies in $D_{Q_\tau}$ for all small $\tau$. So $N_\tau$
increases to $N_0$ as $\tau$ decreases to $0$, and, by dominated convergence with
dominating functions $1$ and $\|h\|$ (M1), $\mu(N_0\setminus N_\tau)\to0$ and $\|m(N_0)-
m(N_\tau)\|\le\int_{N_0\setminus N_\tau}\|h\|\,d\mu\to0$.

The gained set is non-null. If $N_0\setminus X$ were null, then for $k\in K'$ we would
have $k(m(N_0)-m(X))=-\int_{X\setminus N_0}k(h)\,d\mu\le0$, since $k(h)>0$ on $U$; this
contradicts $k(m(N_0)-m(X))\ge v>0$. Fix $\tau$ with $\|m(N_0)-m(N_\tau)\|\le v/2$ and
$\mu(N_0\setminus N_\tau)<\mu(N_0\setminus X)$, and put $N\coloneqq N_\tau$. Since states
have norm $1$, $q(m(N)-m(X))\ge v-v/2=v/2$ for every $q\in S$, which is (i); and
$\mu(N\setminus X)\ge\mu(N_0\setminus X)-\mu(N_0\setminus N)>0$, which is (iii).

\emph{Step D: hull invariance.} We first show: if $q\in S$, $c\ge0$ and $q(h)\ge-c\,P(h)$
for a.e.\ $h\in X$, then $q\in Q_X$. For $q'\in Q_X$ put $T'(q')\coloneqq(q+cq')/(1+c)\in
S$. For a.e.\ $h\in X$, $(q+cq')(h)\ge q(h)+c\,P(h)\ge0$, so $T'$ maps $Q_X$ into itself.
Its iterates satisfy $T'^k(q')-q=(c/(1+c))^k(q'-q)$, so they converge to $q$ in the norm
of $E^\ast$, hence weak-$\ast$; and $Q_X$ is nonempty and weak-$\ast$ closed (Step A).
So $q\in Q_X$.

Now let $q\in Q_{X\setminus N}$ and $c\coloneqq(1-\tau)/\tau$. For $h\in X\cap N$,
$P_{Q_\tau}(h)>0$ gives $P_S(h)>-c\,P(h)$, and $q(h)\ge P_S(h)$. For a.e.\ $h\in
X\setminus N$, $q(h)\ge0$, and $P(h)\ge0$ by Step A, so $q(h)\ge-c\,P(h)$. Hence $q\in
Q_X$. Since $Q_X\subseteq Q_{X\setminus N}$ trivially, $Q_{X\setminus N}=Q_X$ and
$\hull(X\setminus N)=\hull(X)=G\supseteq N_0\supseteq N$. This is (ii). Finally
$\mu\bigl((N\setminus X)\cap\hull(X\setminus N)\bigr)=\mu(N\setminus X)>0$, so $r(N,X)$
is safe by Lemma~\ref{lem:inf-safe}.
\end{proof}

\subsection{The ranking}\label{subsec:inf-ranking}

For $X\in\M_U$ and $\omega\in\Omega$ put $m_\omega(X)\coloneqq m(X)(\omega)=\int_X
h(\omega)\,d\mu$, finite by (M1), and $m_\omega(r)\coloneqq m_\omega(A)-m_\omega(B)$ for a
record $r=(A,B)$. A ranking is a family $(\preceq_\omega)_{\omega\in\Omega}$ of relations
on $\M_U$; dominance, the conditions (TC), (FCC), (Add) and the requirements (G), (IB),
(GC), (C), (ID), (W) are defined as in Section~\ref{sec:framework}, with $h(\omega)$ in
place of $h_i$ and $\omega\in\Omega$ in place of $i$. The record $(A,B)$ is
\emph{guidance at $\omega$} if $A\subseteq\{h(\omega)>0\}$ and $B\subseteq\{h(\omega)<0\}$
a.e.; $G_\omega$ is the set of such records. The \emph{safe-improvement ranking} on $\M_U$
is
\[
  X\preceq_\omega Y\iff r(X,Y)\in G_\omega\cup\{r:r\text{ is safe and }m_\omega(r)>0\}.
\]

\begin{lemma}\label{lem:normingE}
Assume (M2).
\begin{enumerate}[label=(\alph*)]
  \item If $r\in G_\omega$, then $m_\omega(r)\ge0$, with equality if and only if $r$ is
  null.
  \item There is a countable set $\Omega_0\subseteq\Omega$ with $\inf_{\Omega_0}h=
  \inf_\Omega h$ and $\sup_{\Omega_0}h=\sup_\Omega h$ for every $h\in E$.
  \item Let $\Omega'\subseteq\Omega$ satisfy $\inf_{\Omega'}h=\inf_\Omega h$ for every
  $h\in E$. A record that is guidance at every $\omega\in\Omega'$ is a.e.\ of the form
  $(A,\emptyset)$ with $A\subseteq D_S$; it is then guidance at every $\omega\in\Omega$,
  and, if it is non-null, it is safe.
\end{enumerate}
\end{lemma}

\begin{proof}
(a) $m_\omega(r)=\int_A|h(\omega)|\,d\mu+\int_B|h(\omega)|\,d\mu$, which vanishes only if
$A\cup B\subseteq\{h(\omega)=0\}$ a.e.; this set is null by (M2) with $Q=\{\delta_\omega
\}$.
(b) Let $(h_j)$ be dense in $E$, and choose $\omega_{j,l}$ with $h_j(\omega_{j,l})<\inf_\Omega
h_j+1/l$. Let $\Omega_1\coloneqq\{\omega_{j,l}\}$. If $\|h-h_j\|<\eta$, then $\inf_{\Omega_1}
h\le\inf_{\Omega_1}h_j+\eta\le\inf_\Omega h_j+\eta\le\inf_\Omega h+2\eta$; so
$\inf_{\Omega_1}h=\inf_\Omega h$. Applying the same construction to $(-h_j)$ gives a
countable $\Omega_2$ for suprema; put $\Omega_0\coloneqq\Omega_1\cup\Omega_2$.
(c) Let $(A,B)$ be guidance at every $\omega\in\Omega'$. Since $-h\in E$ for $h\in E$,
also $\sup_{\Omega'}h=\sup_\Omega h$. Choosing the points $\omega_{j,l}$ of (b) within
$\Omega'$ gives a countable $\Omega''\subseteq\Omega'$ with $\inf_{\Omega''}h=\inf_\Omega h$
and $\sup_{\Omega''}h=\sup_\Omega h$ for all $h\in E$. Off a countable union of
null sets, every $h\in A$ has $h>0$ on $\Omega''$, so $P_S(h)=\inf_\Omega h\ge0$; removing
the null set $\{P_S=0\}$ (M2), $A\subseteq D_S$ a.e. Likewise every $h\in B$ off a null set
has $\sup_\Omega h\le0$, so $k(h)\le0$ for $k\in K'$, and $h\notin U$; so $B$ is null.
Every $h\in D_S$ satisfies $h(\omega)\ge\inf_\Omega h>0$ for all $\omega$, so the record is
guidance at every world. If $A$ is non-null, the record is safe by
Lemma~\ref{lem:inf-safe}, since $\mu(A\cap D_S)=\mu(A)>0$.
\end{proof}

Part (c) is the only place where the size of $\Omega$ interacts with null sets. Strict
guidance at a single world needs only (a). Example~\ref{ex:PS} shows that the null set
$\{P_S=0\}$ in (c) is essential.

\begin{theorem}\label{thm:inf-main}
Assume (M1) and (M2), and consider the safe-improvement ranking on $\M_U$.
\begin{enumerate}[label=(\alph*)]
  \item Each $\preceq_\omega$ is reflexive, ties a.e.-equal models, and satisfies (G),
  (IB), (TC), (FCC) and (Add). If $X\preceq_\omega Y$ and $X\ne Y$ a.e., then
  $m_\omega(X)>m_\omega(Y)$; so there are no ties between models that differ on a
  non-null set.
  \item For $X,Y\in\M_U$, the following are equivalent: $Y$ weakly dominates $X$ and $Y\ne
  X$ a.e.; $Y$ dominates $X$; $Y$ strictly dominates $X$; $r(Y,X)$ is safe and
  $m_\omega(Y)>m_\omega(X)$ for every $\omega\in\Omega$.
  \item No member of $\CC_U$ has a weak dominator other than itself.
  \item Every $X\in\M_U\setminus\CC_U$ is strictly dominated by some $N\in\CC_U$ with
  $\inf_{q\in S}q(m(N)-m(X))>0$, $N\subseteq\hull(X)=\hull(X\setminus N)$ and
  $\mu(N\setminus X)>0$.
  \item The undominated models are exactly the members of $\CC_U$, and (GC), (ID), (W)
  and (C) hold, the last in the strong form: a model tied at some world with a member of
  $\CC_U$ equals it. If (M3) holds, $\CC_U$ is the class of all coherent models on $U$.
\end{enumerate}
\end{theorem}

\begin{proof}
Apply Lemma~\ref{lem:abstract} with $\M=\M_U$, $\CC=\CC_U$, world set $\Omega$,
$\sigma_\omega=m_\omega$, $d_\omega(r)=m_\omega(r)$, $\mathcal G_\omega=G_\omega$, and
$\mathcal T$ the triples with $(X\triangle Y)\cap(Y\triangle Z)$ null. By
Lemma~\ref{lem:record}(a), the abstract notion of safety is the present one. (A1) is
Lemma~\ref{lem:normingE}(a). (A2): if $r(Y,C)=(Y\setminus C,C\setminus Y)$ is guidance at
every world and $C\in\CC_U$, Lemma~\ref{lem:normingE}(c) with $\Omega'=\Omega$ shows that
$C\setminus Y$ is null and $Y\setminus C\subseteq D_S\subseteq C$ a.e.; since $Y\setminus C$
misses $C$, it is null. (A3) is Theorem~\ref{thm:Hinf}. (A4) follows from
Lemma~\ref{lem:record}(b), the closure of $G_\omega$ under $\sqcup$, and
Lemma~\ref{lem:safe}(c). (A5) is Lemma~\ref{lem:normingE}(c). Lemma~\ref{lem:abstract}(a)
gives (a), apart from (IB), which is Lemma~\ref{lem:record}(c), and (G), which follows from
Lemma~\ref{lem:abstract}(a) because a non-null record satisfying the hypotheses of (G) lies
in $G_\omega\setminus\{0\}$. Lemma~\ref{lem:abstract}(e) gives (Add). (b) follows from
Lemma~\ref{lem:abstract}(a),(f) as in the proof of Proposition~\ref{prop:dominance}. (c)
and (d) are Lemma~\ref{lem:abstract}(b),(c) together with Theorem~\ref{thm:Hinf}. For (e),
Lemma~\ref{lem:abstract}(d) gives the undominated set and strong (C); (GC) holds because
members of $\CC_U$ are coherent; an incoherent model is not in $\CC_U$, so (d) gives (ID);
(W) follows from (c) and (d). The last claim is Lemma~\ref{lem:inf-hull}(d).
\end{proof}

\smallskip\noindent\textbf{Worlds and norming sets.} Let $\Omega_0$ be as in
Lemma~\ref{lem:normingE}(b). The same rule with world set $\Omega_0$ defines a different
dominance relation. For example, let $\Omega=[0,1]$, $E=\operatorname{span}\{\one,\omega,
\omega^2\}\cong\R^3$, $K'$ the Lebesgue mean, $\mu$ the standard Gaussian distribution of
the coefficients, $\Omega_0$ any countable norming set, and $\omega^\ast\in[0,1]\setminus
\Omega_0$. For a small ball $V\subseteq U\setminus\cl D_S$, Lyapunov's theorem shows that
$\{m(A_1)-m(V\setminus A_1):A_1\subseteq V\}=\{\int_Vf h\,d\mu:-1\le f\le1\}$, a convex
symmetric set whose support function $\lambda\mapsto\int_V|\lambda(h)|\,d\mu$ is positive,
so it contains a ball about $0$. Adjoining to $A_1$ a small ball $A_0\subseteq D_S$ disjoint
from $V$ therefore gives a record $r=(A_0\cup A_1,V\setminus A_1)$ with $m(r)=\beta(\omega-
\omega^\ast)^2$ for small $\beta>0$, which is safe by Lemma~\ref{lem:inf-safe} because
$A_0\subseteq D_S\subseteq\hull(V\setminus A_1)$. The models $Y\coloneqq A_0\cup A_1$ and
$X\coloneqq V\setminus A_1$ have $r(Y,X)=r$; so $Y$ strictly dominates $X$ with world set
$\Omega_0$, but not with world set $\Omega$, since $m_{\omega^\ast}(r)=0$ and $r$ is
non-null, so $r$ is not accepted at $\omega^\ast$ (Lemma~\ref{lem:normingE}(a)). The
undominated set is nevertheless the same. Indeed Lemma~\ref{lem:abstract}
applies verbatim with world set $\Omega_0$: (A2) and (A5) hold by
Lemma~\ref{lem:normingE}(c) with $\Omega'=\Omega_0$, and (A3) holds because the margin of
Theorem~\ref{thm:Hinf} is positive at every world. So the undominated models are again
exactly the members of $\CC_U$, and the same holds for every world set $\Omega'$ with
$\Omega_0\subseteq\Omega'\subseteq\Omega$.

\begin{remark}\label{rem:finite-special}
The finite setting is the special case $\Omega=\{1,\dots,n\}$, $E=\R^n$: the states are
the points of $\Delta$, the credal sets agree, and a measure satisfying (M) satisfies (M1),
(M2) (Lemma~\ref{lem:D}(b)) and (M3). Theorem~\ref{thm:Hinf} then gives
Theorem~\ref{thm:H} without Lyapunov's theorem and with any finite Borel measure satisfying
(M1) and (M2), including measures carried by countably many lines $g_k+\R\one$, which are
singular with respect to Lebesgue measure (Lemma~\ref{lem:measure} with $E=\R^n$).
\end{remark}

\begin{remark}\label{rem:orderunit}
The assumption $\one\in E$ can be replaced by: $E$ contains some $e$ with $\inf_\Omega e>
0$. States are then the positive functionals with $q(e)=1$. The map $h\mapsto h/e$ is a
linear bijection of $E$ onto $\tilde E\coloneqq\{h/e:h\in E\}$, bounded with bounded
inverse, so $\tilde E$ is a closed separable subspace of $\ell^\infty(\Omega)$ containing
$\one$. It preserves the sign of every gamble at every world, so it preserves guidance;
$q\mapsto\tilde q$, $\tilde q(h/e)\coloneqq q(h)$, maps states to states and gives
$P_{\tilde Q}(h/e)=P_Q(h)$, so it preserves the sets $D_Q$, (M1)--(M3) for the image
measure, safety and $\CC_U$; and $m_\omega(r)=e(\omega)\tilde m_\omega(\tilde r)$ with
$e(\omega)>0$, so it preserves the ranking. All results of this section transfer, with
margin $(\inf_\Omega e)\tilde v$ in place of $\tilde v$. Some condition of this kind is
needed: if $h(\omega_1)=0$ for every $h\in E$ at some world $\omega_1$, then no comparison at
$\omega_1$ is strict under the rule above, so no model is dominated, and any coherence
notion that some model fails to satisfy violates (GC). Whether a strictly positive gamble
without a positive lower bound suffices (for example $E=c_0$ with $e(n)=1/n$) is open; the
uniform conclusion of Theorem~\ref{thm:Hinf}(i) is impossible there, since every element of
$c_0$ has infimum at most $0$.
\end{remark}

\subsection{Local freedom}\label{subsec:inf-LF}

In this setting, (LF) at $D_K$ is read with the norm topology: there is a norm-open
$N\supseteq\partial D_K\setminus\{0\}$, where $\partial$ is the norm boundary, such that
every Borel $X$ with $X\triangle D_K\subseteq N$ a.e.\ satisfies $X\subseteq U$ a.e. We
call a domain $V\subseteq E$ \emph{one-sided} if there is no nonempty norm-open $O$ with
$O\subseteq V$ a.e.\ and $-O\subseteq V$ a.e. This notion depends only on the a.e.-class
of $V$. In finite dimensions with (M) it is implied by the definition of
Section~\ref{sec:framework}, because Lebesgue-null sets are symmetric: if $O\subseteq V$ and
$-O\subseteq V$ a.e., then $O\subseteq V\cap(-V)$ a.e. It is strictly weaker: if $F$ is a
closed, symmetric, nowhere dense set of positive Lebesgue measure, then $V\coloneqq\Rpp\cup
F$ is not one-sided in the sense of Section~\ref{sec:framework}, but no nonempty open $O$
has $O\subseteq V$ and $-O\subseteq V$ a.e.: since $F$ is closed and nowhere dense, the
open set $O\setminus\cl\Rpp$ would contain a non-null open set outside $V$ unless it is
empty, so $O\subseteq\cl\Rpp$, likewise $-O\subseteq\cl\Rpp$, and $O\subseteq\{0\}$. So Theorem~\ref{thm:inf-cover}(c) extends Lemma~\ref{lem:LF}(c). The
definition of Section~\ref{sec:framework} cannot be used here, because it is not invariant
under a.e.\ modification of $V$ when the null sets of $\mu$ are not symmetric: for $\mu$ as
in Lemma~\ref{lem:measure} with atomic $\nu$, chosen so that no $-g_k$ is a $g_j$, the union
$L$ of the lines $g_k+\R\one$ satisfies $L\cap(-L)=\emptyset$, so $L$ is one-sided in the
sense of Section~\ref{sec:framework}, yet $L$ is a.e.\ equal to $E$. Under (M3),
$U=D_{K'}$ is one-sided: if $O\subseteq U$ a.e.\ and $-O\subseteq U$ a.e., then
Lemma~\ref{lem:inf-hull}(d) gives $O\subseteq U\cap(-U)=\emptyset$. For a credal set $K$
put $K_\perp\coloneqq\{h\in E:k(h)=0\text{ for all }k\in K\}$.

\begin{proposition}\label{prop:inf-LF}
Assume (M2) and (M3), and let $K$ be a credal set. Then $D_K\in\M_U$ and (LF) holds at
$D_K$ if and only if $K'\subseteq K$ and every $k\in K'$ satisfies $k(h)>0$ for all $h\in
C_K\setminus\{0\}$.
\end{proposition}

\begin{proof}
Suppose the condition holds. Then $D_K\subseteq D_{K'}=U$. By
Lemma~\ref{lem:envelopes}(b), $\partial D_K\subseteq\cl D_K\subseteq C_K$, so each $h\in
\partial D_K\setminus\{0\}$ has $k(h)>0$ for every $k\in K'$, hence $P_{K'}(h)>0$ because
the minimum over $K'$ is attained; that is, $h\in U$. Take $N\coloneqq U$. If $X\triangle
D_K\subseteq U$ a.e., then $X\subseteq D_K\cup U=U$ a.e.

Conversely, suppose $D_K\in\M_U$ and (LF) holds with the open set $N$. By
Lemma~\ref{lem:inf-hull}(d), $K'\subseteq K$. The model $D_K\cup N$ satisfies $(D_K\cup N)
\triangle D_K\subseteq N$, so $N\subseteq U$ a.e., and $N\subseteq U$ by
Lemma~\ref{lem:inf-hull}(d). Let $k\in K'$ and $h\in C_K\setminus\{0\}$ with $k(h)=0$.
Then $0\le P_K(h)\le k(h)=0$, so $h\notin D_K$ while $h+t\one\in D_K$ for $t>0$; hence $h\in
\partial D_K\setminus\{0\}\subseteq N\subseteq U$, contradicting $k(h)=0$.
\end{proof}

In finite dimensions the condition of Proposition~\ref{prop:inf-LF} holds exactly when $K$
is full-dimensional and $K'\subseteq\relint K$, by Lemma~\ref{lem:LF}. In infinite
dimensions ``full-dimensional'' and ``relative interior'' are not the right notions: in
the contamination example below, when $\Omega$ is infinite, $p$ is not an algebraic
relative interior point of $K$, yet the condition holds. (For $k=(1-\varepsilon)p+
\varepsilon s\in K$, $p+t(p-k)=(1-\varepsilon)p+\varepsilon((1+t)p-ts)$ lies in $K$ only if
$s\le\frac{1+t}tp$. If $p\{x\}=0$ for some $x$, take $s=\delta_x$; otherwise $\Omega$ is
countable with atoms $p\{x_n\}\to0$, and $s=\sum_k2^{-k}\delta_{x_{n_k}}$ with $p\{x_{n_k}\}
\le4^{-k}$ is not bounded by any multiple of $p$.)

\begin{theorem}\label{thm:inf-cover}
Assume (M1)--(M3), and let $K$ be a credal set.
\begin{enumerate}[label=(\alph*)]
  \item The safe-improvement ranking on $\M_U$ covers $K$ if and only if $K'\subseteq K$
  and every $k\in K'$ is strictly positive on $C_K\setminus\{0\}$. Every ranking on $\M_U$
  that covers $K$ requires this condition. So the safe-improvement ranking is a universal
  cover on $U$.
  \item If $K_\perp=\{0\}$, then $K$ is covered on some domain of the form $D_{K'}$, namely
  with $K'=\{k_0\}$, $k_0\coloneqq\sum_j2^{-j}k_j$ for a weak-$\ast$ dense sequence
  $(k_j)_{j\ge1}$ in $K$.
  \item If $K_\perp\ne\{0\}$, then (LF) fails at $D_K$ on every one-sided domain, so no
  ranking on any one-sided domain covers $K$. If $E$ is infinite-dimensional, this applies
  to every $K$ that is the convex hull of finitely many states, and in particular to every
  precise $K=\{p\}$.
\end{enumerate}
\end{theorem}

\begin{proof}
(a) By Theorem~\ref{thm:inf-main}, requirements (i), (ii), (iii) and (v) of
Definition~\ref{def:cover} hold; $D_K\in\M_U$ iff $K'\subseteq K$, by
Lemma~\ref{lem:inf-hull}(d), and then $D_K\in\CC_U$ is undominated. Requirement (vi) and
$D_K\in\M_U$ are characterised by Proposition~\ref{prop:inf-LF}, whatever the ranking.

(b) $K$ is a compact metrizable space, so a dense sequence exists. The partial sums
$s_J\coloneqq\sum_{j\le J}2^{-j}k_j$ converge in norm, and $s_J/(1-2^{-J})\in K$; since $K$
is weak-$\ast$ closed, hence norm-closed, $k_0\in K$. Let $h\in C_K$ with $k_0(h)=0$. Each
$k_j(h)\ge0$, so $k_j(h)=0$ for all $j$, and by density $k(h)=0$ for all $k\in K$; so $h\in
K_\perp=\{0\}$. Now apply (a) with $K'=\{k_0\}$.

(c) Let $0\ne h\in K_\perp$. Then $P_K(\pm h)=0$ and $\pm h+t\one\in D_K$ for $t>0$, so
$\pm h\in\partial D_K\setminus\{0\}$. If (LF) held on a domain $V$ with the open set $N$,
then, as in the proof of Proposition~\ref{prop:inf-LF}, $N\subseteq V$ a.e.; choose $\rho>0$
with $O\coloneqq\operatorname{Ball}(h,\rho)\subseteq N$ and $-O=\operatorname{Ball}(-h,
\rho)\subseteq N$. Then $O\subseteq V$ a.e.\ and $-O\subseteq V$ a.e., so $V$ is not
one-sided. If $K=\conv\{k_1,\dots,k_m\}$, then $K_\perp=\bigcap_i\ker k_i$ has codimension
at most $m$, so it is nonzero when $E$ is infinite-dimensional.
\end{proof}

Local freedom here is the literal condition, with a norm-open neighbourhood. In finite
dimensions it is automatically uniform; on $E=C(\Omega)$ with $\Omega$ infinite the uniform
version fails (Proposition~\ref{prop:tube}), although in other infinite-dimensional spaces
it can hold. Say that (LF) holds \emph{uniformly} at $D_K$ on $U$ if there is $\delta>0$ such
that the open ball of radius $\delta\|h\|$ about every $h\in\partial D_K\setminus\{0\}$
lies in $U$.

\begin{proposition}\label{prop:tube}
Let $\Omega$ be an infinite compact metric space, $E=C(\Omega)$, and $K'\subseteq K$
credal sets. Then (LF) does not hold uniformly at $D_K$ on $U=D_{K'}$.
\end{proposition}

\begin{proof}
Fix $k\in K'$, a regular Borel probability measure on $\Omega$. Since $\Omega$ is infinite
and metric, it contains infinitely many pairwise disjoint nonempty open sets $V_j$, and
$\sum_jk(V_j)\le1$, so $k(V_j)\to0$. By Urysohn's lemma there are continuous $b_j:\Omega
\to[0,1]$ vanishing outside $V_j$ with $\max b_j=1$; then $k(b_j)\le k(V_j)$. Put $h_j
\coloneqq b_j-P_K(b_j)\one$. Since $0=P_S(b_j)\le P_K(b_j)\le k(b_j)$, we have $P_K(h_j)=0$,
$\|h_j\|\ge1-k(b_j)$, and $h_j\ne0$ because $b_j$ is not constant. So $h_j\notin D_K$,
while $h_j+t\one\in D_K$ for $t>0$; hence $h_j\in\partial D_K\setminus\{0\}$. The gamble $h_j-k(h_j)\one$ is not in $U$, since $k$ vanishes on it,
and its distance from $h_j$ is $k(h_j)\le k(b_j)$. For any $\delta>0$ and large $j$,
$k(b_j)<\delta(1-k(b_j))\le\delta\|h_j\|$, so the ball of radius $\delta\|h_j\|$ about
$h_j$ is not contained in $U$.
\end{proof}

\subsection{Contamination targets}

Let $\Omega$ be a compact metric space and $E=C(\Omega)$, which is separable and contains
$\one$; the states are the regular Borel probability measures on $\Omega$, and $P_S(h)=
\min_\Omega h$. Let $p$ be a Borel probability measure on $\Omega$, $0<\varepsilon\le1$,
and
\[
  K\coloneqq(1-\varepsilon)\{p\}+\varepsilon S,\qquad P_K(h)=(1-\varepsilon)p(h)+
  \varepsilon\min_\Omega h ,
\]
a credal set by Lemma~\ref{lem:envelopes}(c).

\begin{proposition}\label{prop:contamination}
Let $\mu$ satisfy (M1)--(M3), for instance as in Lemma~\ref{lem:measure}.
\begin{enumerate}[label=(\alph*)]
  \item If $p$ has full support, then the safe-improvement ranking on $\M_U$ with
  $U\coloneqq D_{\{p\}}=\{h:p(h)>0\}$ covers $K$. Here $D_K$ is undominated by
  protection (Theorem~\ref{thm:inf-main}(c)), and (LF) holds in the literal form with
  $N=U$.
  \item If $p$ does not have full support, then (LF) fails at $D_K$ on $D_{\{p\}}$.
  \item For every $p$, let $r$ be a Borel probability measure on $\Omega$ with full support
  and $k'\coloneqq(1-\varepsilon)p+\varepsilon r$. Then $k'\in K$, and the safe-improvement
  ranking on $\M_{D_{\{k'\}}}$ covers $K$.
\end{enumerate}
\end{proposition}

\begin{proof}
(a) $p=(1-\varepsilon)p+\varepsilon p\in K$. Let $h\in C_K\setminus\{0\}$. Since $\min h\le
p(h)$, $p(h)\ge(1-\varepsilon)p(h)+\varepsilon\min h=P_K(h)\ge0$. If $p(h)=0$, then
$\varepsilon\min h\ge0$, so $h\ge0$; a nonnegative continuous function with $p(h)=0$
vanishes on the support of $p$, which is $\Omega$, so $h=0$, a contradiction. So $p(h)>0$,
and Theorem~\ref{thm:inf-cover}(a) applies with $K'=\{p\}$; its proof, through
Proposition~\ref{prop:inf-LF}, uses $N=U$.

(b) The open set $\Omega\setminus\operatorname{supp}p$ is nonempty, and by Urysohn's lemma
it carries a continuous $h\ge0$, $h\ne0$, vanishing on the nonempty set $\operatorname{supp}
p$. Then $p(h)=0$ and $\min h=0$, so $h\in C_K\setminus\{0\}$ with $p(h)=0$;
Proposition~\ref{prop:inf-LF} shows that (LF) fails.

(c) $k'\in K$ because $r\in S$. Let $g\in C_K\setminus\{0\}$. Then $k'(g)=P_K(g)+
\varepsilon(r(g)-\min g)\ge\varepsilon(r(g)-\min g)\ge0$. If $k'(g)=0$, then $r(g-\min
g)=0$; since $g-\min g\ge0$ is continuous and $r$ has full support, $g$ is constant, $g=c
\one$, and $c=k'(g)=0$, a contradiction. So $k'$ is strictly positive on $C_K\setminus
\{0\}$, and Theorem~\ref{thm:inf-cover}(a) applies.
\end{proof}

By Theorem~\ref{thm:inf-cover}(c), the precise target $\{p\}$ ($\varepsilon=0$) is not
covered on any one-sided domain when $\Omega$ is infinite, while every contamination
target with $\varepsilon>0$ is covered on a suitable one-sided domain.

\subsection{Coherent sets of desirable gambles}

Call $D\subseteq E$ a \emph{coherent set of desirable gambles} if it is a convex cone,
$0\notin D$, and $D$ contains every $h$ with $\inf_\Omega h>0$.

\begin{proposition}\label{prop:cdg}
Assume (M2). For every coherent set of desirable gambles $D$, the set $Q_D\coloneqq\{q\in S:
q\ge0\text{ on }D\}$ is a credal set and $D_{Q_D}\subseteq D\subseteq C_{Q_D}$. So $D$ is
measurable for the completion of $\mu$, and $D=D_{Q_D}$ a.e.
\end{proposition}

\begin{proof}
(1) $D\supseteq D_S$, which is open and contains $\one$; so the interior $\inter D$ is a
nonempty open convex cone not containing $0$, and $D\subseteq\cl\inter D$, since a convex set
with nonempty interior lies in the closure of its interior.

(2) Let $\psi\in E^\ast$ be nonzero and nonnegative on $D$. For $h\ge0$ and $t>0$,
$h+t\one\in D_S\subseteq D$, so $\psi(h)\ge-t\psi(\one)$; hence $\psi$ is positive. Then
$|\psi(h)|\le\|h\|\psi(\one)$, so $\psi(\one)>0$, and $\psi/\psi(\one)\in Q_D$. So the
nonzero functionals nonnegative on $D$ are exactly the positive multiples of members of
$Q_D$.

(3) By the Hahn--Banach theorem there is $\varphi\in E^\ast$ with $\varphi>0$ on $\inter D$.
By (1), $\varphi\ge0$ on $D$, so $\varphi/\varphi(\one)\in Q_D$ by (2), and $Q_D\ne
\emptyset$. $Q_D$ is weak-$\ast$ closed and convex, so it is a credal set.

(4) By (2) and the bipolar theorem, $\cl D=C_{Q_D}$. Hence $D\subseteq C_{Q_D}$, and
$D_{Q_D}\subseteq\inter C_{Q_D}=\inter\cl D=\inter D\subseteq D$, where $\inter\cl D=
\inter D$ because $D$ is convex with nonempty interior.

(5) $D_{Q_D}$ is open, and $D\setminus D_{Q_D}\subseteq C_{Q_D}\setminus D_{Q_D}=\{P_{Q_D}
=0\}$, which is null by (M2). So $D$ differs from a Borel set by a subset of a null set.
\end{proof}

So, modulo $\mu$-null sets, the coherent models are exactly the coherent sets of desirable
gambles: boundary choices, including those that accept gambles $h$ with $\inf_\Omega h=0$,
are invisible. The same holds in the finite setting, with the same proof. In
Example~\ref{ex:PS}, where (M2) fails, the coherent set of desirable gambles $C\cup R$
differs from every $D_Q$ on a non-null set: if $q_0\in Q$, then $D_Q$ misses $R$; otherwise
$Q\subseteq\{q_s:s\ge s_1\}$ for some $s_1>0$, and $D_Q$ contains the nonempty open set
$\{a<0,\ a+bs_1>0\}$, which misses $C\cup R\subseteq\{a\ge0\}$.

\section{Further properties of the ranking}\label{sec:further}

\subsection{Transitivity}

Transitivity,
\begin{itemize}
  \item[(Tr)] if $X\preceq_iY$ and $Y\preceq_iZ$, then $X\preceq_iZ$,
\end{itemize}
implies (Add), but the safe-improvement ranking does not satisfy it in general.

\begin{example}\label{ex:tr}
Let $n=2$, $K'\coloneqq\{(\tfrac12,\tfrac12)\}$, so that $U=\{h:h_1+h_2>0\}$, and let
$\mu$ be any measure satisfying (M). The model $C\coloneqq U=D_{K'}$ is coherent. Let
$V_a,V_b,V_c$ be the open balls of radius $\tfrac1{10}$ about $a\coloneqq(-1,3)$,
$b\coloneqq(-2,3)$ and $c\coloneqq(1,1)$. They are pairwise disjoint and contained in $U$;
moreover $h_1\in(-1.1,-0.9)$ on $V_a$, $h_1<-1.9$ on $V_b$, and $h_1\in(0.9,1.1)$ on $V_c$.
Put $\delta\coloneqq\min\{\mu(V_a),\mu(V_b),\mu(V_c)\}>0$. Since spheres are null, the
measure of a ball about a fixed centre is continuous in its radius, so there are
concentric balls $A\subseteq V_a$, $B_-\subseteq V_b$ and $B_+\subseteq V_c$ with
$\mu(A)=\delta/2$, $\mu(B_-)=\delta$ and $\mu(B_+)=\delta/4$. Put $B\coloneqq B_-\cup B_+$
and
\[
  Y\coloneqq C\setminus A,\qquad X\coloneqq C\setminus B .
\]
\begin{itemize}
  \item $r(Y,C)=(\emptyset,A)$, and $A\subseteq\{h_1<0\}$, so $r(Y,C)\in G_1$ and
  $Y\preceq_1C$.
  \item $r(X,Y)=(A,B)$. First, $A\subseteq\hull(B)$. Let $q\in Q_B$. Since $B_-$ is not
  null, some $b'\in B_-$ has $q\cdot b'\ge0$. For $a'\in A$, the vector $a'-\tfrac35b'$
  has first coordinate greater than $-1.1+\tfrac35\cdot1.9=0.04$ and second coordinate
  greater than $2.9-\tfrac35\cdot3.1=1.04$, so $a'-\tfrac35b'\gg0$ and $q\cdot a'>
  \tfrac35\,q\cdot b'\ge0$. So $A\subseteq D_{Q_B}=\hull(B)$, and $(A,B)$ is safe by
  Lemma~\ref{lem:safe}(a). Second, $m_1(A)>-1.1\cdot\tfrac\delta2=-0.55\,\delta$ and
  $m_1(B)<-1.9\,\delta+1.1\cdot\tfrac\delta4=-1.625\,\delta$, so $m_1(r(X,Y))>0$. Hence
  $X\preceq_1Y$.
  \item $r(X,C)=(\emptyset,B)$ applies to $C$, so it is not safe; and it is not in $G_1$,
  since $B_+\subseteq\{h_1>0\}$ is not null. Hence not $X\preceq_1C$.
\end{itemize}
So (Tr) fails at world 1, although $m_1(X)-m_1(C)=-m_1(B)>0$. The two steps overlap on
$A$, so (Add) does not apply. Also $r(C,X)=(B,\emptyset)$ has $m_1<0$ and is not in $G_1$,
so $X$ and $C$ are incomparable at world 1.
\end{example}

The example shows the mechanism of protection at work. Removing $B$ from the coherent
model $C$ improves the Brier score at world 1, but the change is not guidance there, and it
applies to $C$; so it is not accepted at world 1. The same change, split into a safe step
and a guided step, is accepted step by step.

\subsection{Relation to the Brier ranking}

By Proposition~\ref{prop:structure}(b), at each world the safe-improvement ranking is
contained in the Brier ranking, and its nontrivial comparisons are strict Brier
improvements. It is far from complete. Every coherent model on $U$ is undominated in it,
while by Remark~\ref{rem:H}(d) every coherent model on $U$ that is not a.e.\ of the form
$U\cap\{\pi\cdot h>0\}$, $\pi\in\Delta$, is uniformly dominated in Brier loss by a model of
that form. By Proposition~\ref{prop:dominance}, a model dominates another in the
safe-improvement ranking exactly when it uniformly dominates it in Brier loss by a safe
change.

One can also define a lexicographic variant of the ranking, using the score
$(m_i,\sum_jm_j)$ at world $i$; we do not use it.

\section{Scope and limits}\label{sec:scope}

\subsection{Universality is relative to a fixed domain}

Theorem~\ref{thm:main} provides one ranking for each domain $U=D_{K'}$. That ranking makes
every coherent model on $U$ undominated and covers exactly the full-dimensional credal
sets whose relative interiors contain $K'$, which are all the credal sets coverable on $U$
(Corollary~\ref{cor:universal}). For any ranking on $\M_{D_{K'}}$, a cover of a
full-dimensional $K$ requires $K'\subseteq\relint K$ (Lemma~\ref{lem:LF}(a),(b)); so two
full-dimensional credal sets with disjoint relative interiors are not covered on a common
domain of this form. Lower-dimensional credal sets are covered on no one-sided domain,
although their coherent models are undominated whenever they are models on $U$
(Proposition~\ref{prop:lowdim}). Sections~\ref{sec:framework}--\ref{sec:main} concern
finitely many worlds and a measure satisfying (M); Lemma~\ref{lem:abstract} and
Section~\ref{sec:infinite} allow any set of worlds. In Section~\ref{sec:infinite} the same
holds with the condition of Theorem~\ref{thm:inf-cover}(a) in place of ``full-dimensional
with $K'\subseteq\relint K$'': every credal set $K$ with $K_\perp=\{0\}$ is covered on some
domain $D_{K'}$, and no credal set with $K_\perp\ne\{0\}$ is covered on any domain that is
one-sided in the sense of Section~\ref{subsec:inf-LF}.

\subsection{Coherence enters the definition of the ranking}

Whether the safe-improvement ranking accepts a comparison that is not guidance depends on
whether the record of the comparison applies to some coherent model. The definition of the
ranking therefore refers to the class of coherent models. The protection of coherent
models (Corollary~\ref{cor:protect}) follows from this definition together with the fact
that guidance alone cannot weakly dominate a coherent model. The dominance of incoherent
models does not follow from the definition: it requires that every incoherent model has a
safe repair that improves its Brier score at every world, which is
Theorem~\ref{thm:H}. Whether a comparison is accepted at world $i$ depends only on its
record, through guidance at $i$, the Brier increment at $i$, and the safety test, which is
the same at every world.

\subsection{Infinitely many worlds}

The following limits apply to Section~\ref{sec:infinite}.
\begin{itemize}
  \item \emph{Separability is essential.} With $\one\in E$, the state space is
  weak-$\ast$ metrizable exactly when $E$ is separable (if $S$ is metrizable, $C(\hat
  \Omega)$ is separable, and $E$ embeds isometrically in it by Lemma~\ref{lem:states}(c)),
  and the proofs use metrizability and a countable norming set. On a nonseparable normed space there is no finite Borel measure
  with full support: for some $\eta>0$ there is an uncountable family of points at mutual
  distance at least $\eta$, and the disjoint balls of radius $\eta/2$ about them would all
  need positive measure. In particular $E=\ell^\infty(\Omega)$ with $\Omega$ infinite, and
  $E=C(\Omega)$ with $\Omega$ compact but not metrizable, lie outside
  Section~\ref{sec:infinite}; their separable subspaces containing $\one$ lie inside it.
  \item \emph{The constants.} The constant $\one$ can be replaced by any gamble bounded
  below by a positive constant (Remark~\ref{rem:orderunit}). Some condition of this kind is
  needed: a world at which every gamble pays $0$ makes every model undominated. Whether a
  strictly positive gamble without a positive lower bound suffices, as for $c_0$ with
  $e(n)=1/n$, is open.
  \item \emph{The null boundary $\{P_S=0\}$.} Condition (M2) for $Q=S$ makes the gambles
  with $\inf_\Omega h=0$ negligible. Without it, guidance alone can strictly dominate a
  coherent model (Example~\ref{ex:PS}). The same condition makes boundary choices of
  coherent sets of desirable gambles invisible (Proposition~\ref{prop:cdg}).
  \item \emph{Dependence on $\mu$.} The coherent models are determined by their credal
  sets, whatever full-support $\mu$ is used (Lemma~\ref{lem:inf-hull}(d)). But models are
  classes modulo $\mu$-null sets, so which models are incoherent, and hence the ranking,
  depends on $\mu$, and admissible measures can be mutually singular. For instance, if
  $\mu$ and $\mu'$ are built as in Lemma~\ref{lem:measure} from dense sequences $(g_k)$ and
  $(g'_k)$ in $E_0$ with no common term, they are carried by disjoint unions of lines
  $g_k+\R\one$ and $g'_k+\R\one$. If $L'$ is the union of the lines $g'_k+\R\one$, then
  for a credal set $Q\supseteq K'$, $D_Q\setminus L'$ is coherent for $\mu$ (it equals $D_Q$ a.e.) but is $\mu'$-null, and so
  incoherent for $\mu'$, since every coherent model contains the nonempty open set $D_S$.
  \item \emph{Local freedom is literal, not uniform.} (LF) is the condition of
  Section~\ref{sec:framework} with a norm-open neighbourhood. Its uniform version, which
  is automatic in finite dimensions, fails for every target on $C(\Omega)$ with $\Omega$
  infinite (Proposition~\ref{prop:tube}).
  \item \emph{States and countable additivity.} States are positive normalised linear
  functionals on $E$. They need not be countably additive on $\Omega$, but every state is
  given by a countably additive regular Borel probability on the compact metrizable space
  $\hat\Omega$ of Lemma~\ref{lem:states}(c). For $E=C(\Omega)$ with $\Omega$ compact
  metric, $\hat\Omega$ can be identified with $\Omega$.
  \item \emph{Coherent models.} Modulo $\mu$-null sets, the coherent models $D_Q$ are
  exactly the coherent sets of desirable gambles (Proposition~\ref{prop:cdg}).
\end{itemize}

\subsection{Representation and justification}

Theorems~\ref{thm:main} and~\ref{thm:inf-main} show that the coherent models on $U$ can be represented as the
undominated models of a ranking satisfying (G), (IB), (TC), (FCC) and (Add), with every
incoherent model strictly dominated by a coherent one and every dominance relation a Brier
improvement at every world, uniform when $\Omega$ is finite; the repair $N$ of
Theorem~\ref{thm:Hinf} improves uniformly in every case. It does not derive coherence from accuracy considerations that
are specified independently of coherence, since the ranking's definition refers to
coherence through the safety test.

\begin{thebibliography}{9}

\bibitem{KCM}
J.~Konek and C.~Campbell-Moore.
\newblock Covering coherent sets of desirable gambles.
\newblock Draft, 2026.

\bibitem{Lindenstrauss}
J.~Lindenstrauss.
\newblock A short proof of {L}iapounoff's convexity theorem.
\newblock \emph{Journal of Mathematics and Mechanics}, 15:971--972, 1966.

\bibitem{Lyapunov}
A.~Lyapunov.
\newblock Sur les fonctions-vecteurs compl\`etement additives.
\newblock \emph{Izvestiya Akademii Nauk SSSR, Seriya Matematicheskaya}, 4:465--478,
1940.

\bibitem{Rockafellar}
R.~T. Rockafellar.
\newblock \emph{Convex Analysis}.
\newblock Princeton University Press, 1970.

\bibitem{Sion}
M.~Sion.
\newblock On general minimax theorems.
\newblock \emph{Pacific Journal of Mathematics}, 8:171--176, 1958.

\end{thebibliography}

\end{document}
